\chapter{Lectores bidireccionales \texorpdfstring{\(K_D\)}{K(D)} y \texorpdfstring{\(K_\Omega\)}{K(Omega)}: perfiles, coincidencias, torres \texorpdfstring{\(90/120\)}{90/120} y operadores de las familias de trayectorias}
\label{chap:lectores-bidireccionales-torres-operadores}

\section{Lectores \(K_D\) y \(K_\Omega\): perfiles dodecafásicos, dominio de rutas y codominio espectral}
Para cada ruta orientada se definen
\[
 K_D=\left(D_{27}+D_{36},D_1,\frac{D_4-D_3}{2}\right),
 \qquad K_\Omega=(\Omega,54,P_6).
\]
El censo exacto de las 324 rutas contiene catorce perfiles \(K_D\), nueve
perfiles \(K_\Omega\), cuarenta pares y dieciocho coincidencias, todas con
valor \((80,54,6)\). Estas cifras describen lectores, no especies materiales.

\section{Semigrupo y operadores de fibra}
Las alturas globales forman
\[
 \langle90,120\rangle=\{90,120\}\cup\{30r:r\ge6\},
\]
con álgebra \(\Bbbk[X_{90},Y_{120}]/(X_{90}^4-Y_{120}^3)\). Los coeficientes
de torre convergentes refinan el carácter y definen operadores positivos sobre
las trece multisecciones.

\section{Construcción electrónica, transductor de torres y reloj de ruta}
El electrón central proporciona el caso de rango uno; el resto de las
multisecciones exige el operador de fibra, el acoplamiento y la lectura
espectral. El transductor conserva por separado longitud temporal, memoria,
acarreo y altura armónica.

\section{Registro dodecafásico y memoria \texorpdfstring{\(1+5+6\)}{1+5+6} del TPK}

El selector operativo, el calendario y el registro dodecafásico fijan la
descomposición de memoria que utilizan los lectores de masa. La identidad
\(1+5+6\) pertenece al estado temporal del TPK y no se reconstruye desde las
salidas espectrales.

\begingroup
\let\chapter\section
\let\section\subsection
\let\subsection\subsubsection
\let\subsubsection\paragraph
\let\clearpage\relax
\let\cleardoublepage\relax
\section{Sistema temporal dodecafásico orientado del TPK}
\label{sec:ciclo-dodecafasico-tpk-rev7}
\index{TPK!sistema dodecafásico}\index{periodicidad de orden 108}

\HMTClaim{HMT-R-TPK-GAMMA108}

\subsection{Operador de selección de ruta y transporte al observable dimensional}

La función ternaria de fase operativa es
\[
M_{\rm ph}(r)=
\begin{cases}
 +1,&r\in\{1,4,7\},\\
 -1,&r\in\{2,5,8\},\\
 0,&r\in\{3,6,9\}.
\end{cases}
\]
Sus tres valores son las entradas escalares con las que
\(\operatorname{Sel}_t\) activa, respectivamente, el cursor aditivo, el
cursor multiplicativo o el contador neutral. El muestreo estroboscópico es
\[
 \operatorname{Obs}_3(a)=(a_3,a_6,a_9,\ldots).
\]
La selección \(\operatorname{Sel}_t[\,M_{\rm ph}(g(t))\,]\) decide qué
evaluación opera; el transporte cardinal desplaza el estado; y
\(\operatorname{Obs}_3\) selecciona qué instantes se publican.

El transporte cardinal actúa por cuatro traslaciones
\[
 T_d(i,j)=(i,j)+\nu_d\pmod9,
\]
\[
 \nu_N=(-1,0),\qquad \nu_E=(0,1),\qquad
 \nu_S=(1,0),\qquad \nu_O=(0,-1).
\]
Los cursores aditivo y multiplicativo conservan posiciones y orientaciones
separadas. La coordenada TRIT, la función \(M_{\rm ph}\), el operador
\(\operatorname{Sel}_t\), las cuatro traslaciones y
\(\operatorname{Obs}_3\) son, por ello, objetos de tipos distintos.

\subsection{Calendario y registro dodecafásicos}

El núcleo topológico de fase no termina en la actualización de una posición
sobre \(\mathbb Z_9^2\). Su estado temporal mínimo conserva el orden de los
eventos, la orientación, el retorno y la memoria de las transiciones. Para
una órbita de \(108\) pasos se introduce la partición ordenada
\[
 E_m=\{9(m-1)+1,\ldots,9m\},
 \qquad m=1,\ldots,12,
\]
y se denota su soporte ordenado por
\[
 \Gamma_{108}=E_1\sqcup\cdots\sqcup E_{12},
 \qquad 108=12\cdot9.
\]
Equivalentemente,
\[
 \Theta_{108}=\mathbb Z/108\mathbb Z,\qquad
 \beta(\theta)=
 \left[\left\lfloor\frac{\widetilde\theta}{9}\right\rfloor\right]_{12},
 \qquad
 r(\theta)=1+(\theta\bmod9).
\]
La coordenada \(\beta\) selecciona uno de doce sectores y \(r\) localiza la
fase dentro de su ventana nonádica.

\begin{proposition}[Extensión central no escindida del reloj observable]
\label{prop:extension-central-108-no-escindida-rev8}
Escribiendo \(C_n=\mathbb Z/n\mathbb Z\), el calendario observable forma la
sucesión exacta
\[
 0\longrightarrow C_{12}
 \xrightarrow{\ \iota\ }C_{108}
 \xrightarrow{\ p\ }C_9
 \longrightarrow0,
 \qquad
 \iota([b]_{12})=[9b]_{108},\qquad
 p([\theta]_{108})=[\theta]_9.
\]
Para la sección conjuntista \(s([r]_9)=[r]_{108}\), con
\(0\leq r<9\), la costura viene dada por el cociclo de acarreo
\[
 c(r,r')=
 \left[\left\lfloor\frac{r+r'}{9}\right\rfloor\right]_{12},
 \qquad
 s(r)+s(r')=s([r+r']_9)+\iota\bigl(c(r,r')\bigr).
\]
La extensión no se escinde. Para la acción trivial,
\[
 H^2(C_9,C_{12})
 \simeq C_{12}/9C_{12}
 \simeq C_3,
\]
y \(c\) representa la clase generadora \([1]\).
\end{proposition}

\begin{proof}
La imagen de \(\iota\) es el subgrupo de los múltiplos de nueve en
\(C_{108}\), que coincide con \(\ker p\); de aquí la exactitud. La identidad
del cociclo es la división euclídea de \(r+r'\) por nueve. Si la extensión
central se escindiera, el grupo medio sería isomorfo a
\(C_{12}\times C_9\). Sin embargo,
\[
 \exp(C_{108})=108,\qquad
 \exp(C_{12}\times C_9)=\operatorname{mcm}(12,9)=36,
\]
lo que lo impide. La identificación cohomológica para un grupo cíclico con
acción trivial da \(H^2(C_9,C_{12})\simeq C_{12}/9C_{12}\), y el acarreo
unitario proyecta al generador de ese cociente.
\end{proof}

La no escisión tipa el retorno: nueve iteraciones cierran la coordenada en
\(C_9\) y trasladan una unidad en \(C_{12}\); ciento ocho iteraciones cierran ambas
coordenadas de este reloj. El calendario \(C_{108}\) es la proyección
temporal del estado enriquecido \(\mathcal X_{\rm TPK}^{\rm enr}\):
orientación, hoja, frontera, cilindro, genealogía y memoria profunda
permanecen en este último y pueden
distinguir historias con el mismo calendario observable.

Si \(g_0(\theta)=\theta\bmod9\), entonces
\[
 g_0(\theta+9)=g_0(\theta),\qquad
 \beta(\theta+9)=\beta(\theta)+1\pmod{12}.
\]
El testigo
\[
 (g_0,r,\beta,30^\circ\beta)(49)=(4,5,5,150^\circ),
\]
\[
 (g_0,r,\beta,30^\circ\beta)(58)=(4,5,6,180^\circ)
\]
muestra exactamente qué retorna y qué conserva memoria: la fase nonádica y
la posición interna coinciden, mientras el sector dodecafásico avanza uno.
El calibre regional \(\operatorname{diag}_c\) de las regiones de \(\pi\) no
es esta coordenada \(\beta\). Tampoco existe aquí un morfismo
\(30^\circ\beta\to h\in\langle90,120\rangle\); en particular,
\(150\notin\langle90,120\rangle\). El par \(49/58\) es un testigo de acarreo,
no un selector regional ni espectral.

Para una historia enriquecida \(x\) de TPK, sean \(\tau_m\) la subruta
ordenada sobre \(E_m\), \(\sigma_m\) su orientación, \(\kappa_m\) el acarreo
que cruza la frontera del bloque y \(\Lambda_m\) el estado acumulado del
registro. El objeto
\[
 \mathcal R_{12}
 =\bigl((E_m,\tau_m,\sigma_m,\kappa_m,\Lambda_m)\bigr)_{m=1}^{12}
\]
es el sistema temporal dodecafásico orientado de TPK. No constituye una cuarta
teoría añadida a APP, TRIT y TPK: es la forma temporal interna con la que
TPK ordena una órbita, conserva sus retornos y publica un registro
reversible. Se distinguen las proyecciones canónicas
\[
 \mathcal X_{\rm TPK}^{\rm enr}
 \longrightarrow \mathcal R_{12}
 \longrightarrow \Theta_{108}.
\]
La primera conserva \((\tau_m,\sigma_m,\kappa_m,\Lambda_m)\); la segunda
publica sólo la sucesión observable y no puede reconstruir por sí sola el
estado firmado. El registro enriquecido se denota \(\mathcal R_{12}\) y la
rotación de sectores se denota \(C_{12}\).

El cierre observable del calendario satisface
\[
 F_{\rm obs}^{108}=I.
\]
Las posiciones y orientaciones ya se restauran tras \(54\) iteraciones,
mientras \(\mathcal R_{12}\) distingue la posición temporal dentro de los
doce sectores. Aquí \(F_{\rm obs}\) es la permutación observable, mientras
\(\mathcal G_9\) es la poda dentro de la relación de extensión y frontera.
La monodromía es la propiedad conjunta
\[
 g(K+9)=g(K),\qquad v_{K+9}\ne v_K,
\]
no el nombre aislado de la poda. El cierre cronológico devuelve el calendario
observable; el retorno nonádico conserva la fase y modifica prefijo, cilindro
y memoria.

También se distinguen dos palabras de doce componentes:
\[
 U_{12}^{\rm cat}=U_6\Vert U_6,
 \qquad
 U_{12}^{\rm sgn}=\mathcal H_{12}(K).
\]
La primera concatena dos emisiones de longitud seis. La segunda es una
coordenada firmada del estado enriquecido y queda ligada reversiblemente a
\(K\) por la transformación dodecafásica \(\mathcal H_{12}\). Ni sus dominios
ni sus funciones son intercambiables; en particular, la palabra concatenada
no construye por sí sola la coordenada firmada.

\begin{definition}[Registro integral de los doce eventos]
\label{def:registro-dodecafasico-rev7}
Para una ruta de cifras \(d_1,\ldots,d_{108}\), sea
\[
 \mathcal D_{\rm evt}=\{2,4,6,8\},\qquad
 \Sigma_{12}=(+,+,+,-,-,-,+,+,+,-,-,-).
\]
Su palabra orientada de eventos es
\[
 e_m=\Sigma_{12}(m)\,
 \#\{t\in E_{m+1}:d_t\in\mathcal D_{\rm evt}\},
 \qquad 0\le m<12.
\]
Las dos semicoronas se emparejan mediante
\[
 p_i=e_i+e_{i+6},\qquad
 a_i=e_i-e_{i+6},\qquad 0\le i<6.
\]
Se definen el conteo bruto, las cinco diferencias respecto del sexto par y
las seis antisimetrías por
\[
 s_C=\sum_{i=0}^{5}p_i,\qquad
 M_i=p_i-p_5\quad(0\le i<5),\qquad
 A_6=(a_0,\ldots,a_5).
\]
El registro integral es
\[
 \mathscr L_{12}(e)
 =\bigl(s_C,M_0,\ldots,M_4,a_0,\ldots,a_5\bigr).
\]
\end{definition}

\begin{theorem}[Coordenadas exactas de la memoria dodecafásica]
\label{thm:memoria-1-5-6}
El mapa \(e\mapsto\mathscr L_{12}(e)\) es inyectivo sobre su imagen. La
inversa viene dada por
\[
 p_5=\frac{s_C-\sum_{i=0}^{4}M_i}{6},\qquad
 p_i=p_5+M_i\quad(0\le i<5),
\]
\[
 e_i=\frac{p_i+a_i}{2},\qquad
 e_{i+6}=\frac{p_i-a_i}{2}.
\]
En particular, la distribución de coordenadas es
\[
 \boxed{1+5+6=12}.
\]
\end{theorem}

\begin{proof}
Por definición,
\[
 s_C=\sum_{i=0}^{4}(p_5+M_i)+p_5
 =6p_5+\sum_{i=0}^{4}M_i.
\]
Esto determina \(p_5\) y después todos los \(p_i\). Sumar y restar las
ecuaciones \(p_i=e_i+e_{i+6}\) y \(a_i=e_i-e_{i+6}\) recupera los doce
eventos. La divisibilidad por seis y las seis condiciones de paridad
caracterizan la imagen integral; no intervienen una masa ni una unidad.
\end{proof}

La división por seis no se aplica durante el conteo y no es un coeficiente
primitivo. Primero se forma el conteo orientado bruto \(s_C\); después se
reconstruye el registro de doce eventos; sólo entonces aparece
\(s_C/6\) en el carácter. Equivalentemente,
\[
 W_\pm=6n_A\pm s_C
\]
son enteros y la factorización reticular posee forma normal de Smith
\((1,12)\). Esta precedencia será obligatoria en la ley de masas.

Los enteros \(90\) y \(120\) reaparecen más adelante en seis construcciones
de tipos distintos. Para impedir que el nombre del cardinal suplante al
objeto, se fija desde ahora el siguiente convenio:

\begin{enumerate}
 \item \(\mathcal E_{\rm dod}=(D_3,D_4,Q):\mathbb Z^{12}\to\mathbb Z^{25}\),
       extractor integral que reconstruye \(K\);
 \item \(D_{\rm raw}^{90,120}(q)\), diferencia diagnóstica de dos series en
       una carta fijada;
 \item \(\iota_{6,8}:\mathcal F_6\hookrightarrow\mathcal F_8\), inclusión
       combinatoria de cardinalidades \(90\) y \(120\);
 \item \(K_{6\to8}=12\Delta S_{90}-\Delta S_{120}\), funcional angular
       definido después de fijar canales, series y normalización;
 \item \(\mathfrak S=\langle90,120\rangle\), semigrupo infinito de alturas
       espectrales;
 \item \(\varepsilon_{\rm res}^{\circ}\), residuo de una identidad angular
       exacta, no una extracción finita de cualquiera de los cinco objetos
       anteriores.
\end{enumerate}

La coincidencia parcial de índices no establece igualdad de operadores,
dominios, codominios ni fuerza probatoria. En particular,
\(K_{6\to8}\) usa dos subcolas de la jerarquía, pero no es la jerarquía
completa; y ninguna de esas subcolas forma un bloque independiente de la ley
de masas.

\begin{remark}[Separación canónica de los objetos temporales]
La notación distingue el registro orientado \(\mathcal R_{12}\), el
calendario observable \(\Theta_{108}\), la rotación sectorial \(C_{12}\) y
la órbita de transporte \(\Gamma_{108}\), cada uno con su dominio, su
codominio y su ley de composición.
\end{remark}

\endgroup

\begingroup
\let\chapter\section
\let\section\subsection
\let\subsection\subsubsection
\let\subsubsection\paragraph
\let\clearpage\relax
% Vista pública derivada de 73_rev2_electron_lectores_torres.tex: cuerpo exacto y único retítulo académico del lector temporal de 108 estados. El fragmento del handoff permanece byte-exacto.
\section{Construcción completa del electrón beta: ruta central, firma, carácter y masa rectora}
\Needspace{7\baselineskip}
\subsection{Ruta central del electrón y secuencia de registros genealógicos}
El electrón ocupa el único centro \((5,5)\), con orientación activa \(++\):

\[
 \Gamma_e=(5,5,++),\qquad
 \tau_{12}=+++---+++---.
\]
Sus lectores de desplazamiento son

\[
 (D_1,D_3,D_4,D_{27},D_{36})=(54,56,68,48,32).
\]
No deben fundirse:

\begin{enumerate}
\item la firma cruda \((24,0,12,0,-12)\);
\item la firma par CT--108 \((-12,-4,12,-6,12)\);
\item el vector de doce eventos;
\item el origen afín \(v_e=0\) del carácter relativo.
\end{enumerate}
Tampoco debe confundirse \(\Gamma_e\) con la palabra \(w_e=201101\) del lector de la constante de Euler \(e\).

\Needspace{7\baselineskip}
\subsection{Selector basal del electrón y transición al valor beta rector}
El selector interno, construido sin usar la masa observada, cierra

\[
 -22=-\left|((\mathbb Z/9\mathbb Z)\times\mathbb F_3)\setminus\iota(R_\pi)\right|,
 \qquad
 +15=|R_\pi\times\mathbb F_3|.
\]
La lectura de vacancias aporta además

\[
 V_e=\begin{pmatrix}21&22\\6&7\end{pmatrix},\qquad
 \det V_e=15.
\]
El signo negativo de 22 registra déficit/corte; el positivo de 15, retención/área orientada. La matriz es local al electrón y no se extiende como selector universal.

La etapa basal es

\[
 e_0=\frac{\sqrt3}{4}
 \left[\exp\!\left(\frac{\varphi}{\pi^2}\right)
 -22\alpha_{\rm HMT}^3\right](1+15\Delta_4),
\]
\[
 e_0=0.510998922488\ldots\ \mathrm{MeV}.
\]
Las genealogías enteras son

\[
 22=\frac{396}{18}=27-5=24-2,
 \qquad
 15=\frac{270}{18}=\binom62=\binom64.
\]
\Needspace{7\baselineskip}
\subsection{Lectores bidireccionales y beta}
Para toda ruta,

\[
 K_D(\Gamma)=\left(D_{27}+D_{36},D_1,\frac{D_4-D_3}{2}\right),
\]
\[
 K_\Omega(\Gamma)=(\Omega_{90,120},54,P_6).
\]
Sobre las 324 rutas aparecen 14 perfiles \(K_D\), 9 perfiles \(K_\Omega\), 40 pares conjuntos y 18 coincidencias. El único valor común es

\[
 (80,54,6).
\]
Este 80 procede del lector bidireccional de la ruta electrónica, no del período 80 de la palabra de Euler. Por ello

\[
 \kappa_e=80-54\alpha_{\rm HMT}+6\Delta_4,
\]
\[
 \boxed{
 \mathfrak e_e^\beta=e_0R_{\rm act}^{\kappa_e}
 =0.5109989506900008086082571648\ldots\ \mathrm{MeV}.}
\]
Esta es la salida electrónica rectora. El basal no la sustituye. La igualdad con la masa de la antipartícula procede de la conjugación de ruta y de la paridad del operador de masa, no de añadir una segunda fila numérica.

\Needspace{7\baselineskip}
\section{Realización física por acoplamiento espectral}
\Needspace{7\baselineskip}
\subsection{Obstrucción a la selección puntual equivariante de una multisección no central}
Sea \(F\subset\mathcal R_{324}\) una fibra de rutas compatible con una clase estructural. Su espacio lineal es

\[
 \mathcal H_F=\bigoplus_{\Gamma\in F}\mathbb C e_\Gamma.
\]
El carácter central y los lectores definen operadores diagonales:

\[
 \widehat{\mathcal A}_F^{(0)}e_\Gamma
 =\chi_{\rm full}(\sigma_\Gamma,\eta_\Gamma)e_\Gamma,
\]
\[
 \widehat K_F^r e_\Gamma=\kappa_r(\Gamma)e_\Gamma,
 \qquad r\in\{D,\Omega\}.
\]
La corrección two--way es

\[
 \widehat{\mathcal M}_F^{\beta,r}
 =R_{\rm act}^{\widehat K_F^r/2}
  \widehat{\mathcal M}_F^{(0)}
  R_{\rm act}^{\widehat K_F^r/2}.
\]
Cuando \(\dim\mathcal H_F>1\), estos operadores tienen, en general, varios autovalores. Elegir uno porque se aproxima a una masa conocida invertiría la causalidad. Tampoco es correcto exigir que un nombre físico señale una celda única: la especie no central reside en una representación de la fibra.

\Needspace{7\baselineskip}
\subsection{Funtor espectral de la realización física: núcleo positivo de acoplamiento, subespacio persistente y medida espectral}
Una preparación física \(P\) aporta datos de representación, no una masa:

\[
 \mathfrak r(P)=
 (q,\,j,\,\epsilon_{2\pi},\,\epsilon_{4\pi},\,
  \text{estadística},\,\text{generación},\,\text{sabor},\,
  \text{órbita de color},\,\text{canal de interacción}).
\]
Sea \(\mathfrak M\) un acoplamiento compatible con esa representación. Se define un núcleo positivo

\[
 K_{P,\mathfrak M}:F\times F\longrightarrow\mathbb C
\]
que satisface:

\begin{enumerate}
\item \(K=K^*\) y \(K\ge0\);
\item \(K(\Gamma,\Gamma')=0\) si las rutas violan carga, paridad, color, estadística o
canal de \(P\);

\item \(K\) conmuta con las simetrías que no rompe el acoplamiento;
\item \(K(C_M\Gamma,C_M\Gamma')=\overline{K(\Gamma,\Gamma')}\) bajo conjugación;
\item sus entradas se calculan con registro genealógico, frontera, memoria, acarreo y holonomía, sin
magnitudes externas.

\end{enumerate}
El operador de evolución acoplado es

\[
 \mathcal U_{P,\mathfrak M}
 =K_{P,\mathfrak M}^{1/2}\,\mathcal U_{\rm TPK}\,
  K_{P,\mathfrak M}^{1/2}.
\]
El subespacio persistente se obtiene mediante la proyección espectral sobre los modos que sobreviven al retorno y a la poda:

\[
 \Pi_{P,\mathfrak M}^{\rm surv}
 =\mathbf1_{\mathcal S_{\rm surv}}(\mathcal U_{P,\mathfrak M}).
\]
La realización física correcta es el funtor espectral

\[
 \boxed{
 \mathcal S_{\rm phys}^{\rm spec}(P,\mathfrak M)
 =\left(
 \mathcal H_{P,\mathfrak M}^{\rm surv},
 \widehat{\mathcal M}_{P,\mathfrak M},
 \mu_{P,\mathfrak M}
 \right),}
\]
donde

\[
 \mathcal H_{P,\mathfrak M}^{\rm surv}
 =\Pi_{P,\mathfrak M}^{\rm surv}\mathcal H_F,
\]
\[
 \widehat{\mathcal M}_{P,\mathfrak M}
 =\Pi_{P,\mathfrak M}^{\rm surv}
  \widehat{\mathcal M}_F
  \Pi_{P,\mathfrak M}^{\rm surv},
\]
y \(\mu_{P,\mathfrak M}\) es su medida espectral en el estado preparado.

Cuando el acoplamiento se especifica por una representación irreducible \(\lambda\) de un grupo finito \(G_P\), el proyector correspondiente se obtiene sin elegir una ruta distinguida:

\[
 P_{P,\lambda}
 =\frac{\dim\lambda}{|G_P|}
  \sum_{g\in G_P}\overline{\chi_\lambda(g)}\,U(g).
\]
La salida bidireccional es entonces el par comprimido

\[
 \mathbb M(P,\lambda)=
 \left(
 P_{P,\lambda}\widehat{\mathcal M}^{\beta,D}P_{P,\lambda},
 P_{P,\lambda}\widehat{\mathcal M}^{\beta,\Omega}P_{P,\lambda}
 \right).
\]
Si ambos operadores conmutan con la representación irreducible, el lema de Schur fuerza su escalaridad sobre el bloque. Si no conmutan, el resultado correcto es el espectro del bloque mínimo invariante.

\Needspace{7\baselineskip}
\subsection{Naturalidad del funtor espectral bajo entrelazadores de la dinámica TPK, el acoplamiento y el operador de masa}
Si \(V:F\to F'\) entrelaza la dinámica TPK, el núcleo de acoplamiento y los operadores de masa, entonces

\[
 V\Pi_{P,\mathfrak M}^{\rm surv}
 =\Pi_{P',\mathfrak M'}^{\rm surv}V,
 \qquad
 V\widehat{\mathcal M}_{P,\mathfrak M}
 =\widehat{\mathcal M}_{P',\mathfrak M'}V.
\]
Por el cálculo funcional, la medida espectral se transporta por \(V\). La salida no depende de una numeración de las celdas ni de la elección de una base dentro de una órbita. Esta es la condición que debe sustituir a todo selector nominal no equivariante.

\Needspace{7\baselineskip}
\subsection{Condiciones de existencia de una masa escalar bajo el lector dimensional}
Una masa escalar está determinada en cualquiera de los casos siguientes:

\begin{enumerate}
\item el subespacio persistente tiene rango uno;
\item una representación irreducible fija un único proyector espectral;
\item el protocolo de medida define un estado positivo normalizado
\(\omega_{P,\mathfrak M}\) sobre el álgebra de observables;

\item una resonancia define un polo aislado del resolvente.
\end{enumerate}
En el tercer caso,

\[
 m(P,\mathfrak M)
 =\omega_{P,\mathfrak M}
  (\widehat{\mathcal M}_{P,\mathfrak M}).
\]
En ausencia de uno de estos mecanismos, la salida científica es el operador, su espectro y sus multiplicidades. Esa salida no es incompleta: es el objeto que la dinámica ha determinado.

\Needspace{7\baselineskip}
\subsection{Construcción finita del núcleo de masa mediante filtros de representación, persistencia y acoplamiento, seguida de normalización irreducible}
El núcleo se obtiene en cuatro etapas finitas:

\begin{enumerate}
\item \textbf{Filtro de representación.} Se retienen las rutas cuyo sector, rama de carga,
color, holonomía \(2\pi/4\pi\), paridad y estadística coinciden con \(P\).

\item \textbf{Filtro de persistencia.} Se usa el orden parcial formado por retorno de clase,
retorno de órbita, retorno de celda, separación de prefijos, deuda de acarreo y estabilidad decimal. El resultado puede ser una órbita de Pareto; no se fuerza unicidad.

\item \textbf{Filtro de acoplamiento.} Se evalúa la incidencia de frontera y el canal de
interacción. Rutas desacopladas reciben peso cero.

\item \textbf{Normalización.} Sobre cada componente irreducible se normaliza la traza del
núcleo a uno.

\end{enumerate}
El atlas orientado proporciona todos los datos de las dos primeras etapas. La tercera exige que cada realización física declare su canal; la cuarta es canónica una vez fijadas las componentes irreducibles.

\Needspace{7\baselineskip}
\section{Transductor genealógico de torres}
\Needspace{7\baselineskip}
\subsection{Par de palabras asociadas a una misma ruta genealógica}
Cada \(\Gamma\in\mathcal R_{324}\) determina simultáneamente una palabra temporal

\[
 w_\Gamma\in\mathbb F_3^{108}
\]
y una palabra dodecafásica firmada

\[
 \tau_\Gamma\in\{-1,0,+1\}^{12}.
\]
La primera conserva desplazamientos a lo largo de CT--108; la segunda conserva la orientación acumulada de los doce acontecimientos. Las dos palabras están ligadas por el mismo registro genealógico, pero no son intercambiables. El refinamiento de torres lee ambas y conserva, además, hoja, acarreo, devanado y memoria de retorno.

\Needspace{7\baselineskip}
\subsection{Correspondencia entre altura armónica y longitud temporal}
Escribamos

\[
 \mathfrak S=30\langle3,4\rangle
 =\{90,120\}\cup\{30r:r\ge6\}.
\]
Para \(h=30r\), la longitud temporal asociada es

\[
 \ell_h=9r.
\]
En particular,

\[
 90\longleftrightarrow27,\qquad
 120\longleftrightarrow36,\qquad
 360\longleftrightarrow108.
\]
Esta correspondencia no identifica la torre con el tiempo: entrelaza la altura del operador armónico con la longitud del segmento de ruta que la genera.

\Needspace{7\baselineskip}
\subsection{Lector de desplazamiento del ciclo temporal de 108 estados}
Definimos la discrepancia cíclica

\[
 D_\Gamma(n)=
 \sum_{t=0}^{107}
 \mathbf1_{\{w_{\Gamma,t}\ne
 w_{\Gamma,t+n\!\!\!\pmod{108}}\}}.
\]
El coeficiente relativo de la rama temporal es

\[
 \boxed{\eta_{30r}^{D}(\Gamma)
 =D_\Gamma(9r)-D_{\Gamma_e}(9r).}
\]
La resta no usa una masa como origen: usa la ruta central como origen afín del carácter relativo. Se cumple

\[
 \eta_{90}^{D}+\eta_{120}^{D}
 =[D_\Gamma(27)+D_\Gamma(36)]-80,
\]
de modo que la primera capa prolonga exactamente el primer componente del lector \(K_D\).

\Needspace{7\baselineskip}
\subsection{Lector de acumulación dodecafásica}
Para \(r\ge1\), sea

\[
 A_\Gamma(r)=
 \sum_{j=0}^{11}
 \left(
 \sum_{u=0}^{r-1}
 \tau_{\Gamma,j+u\!\!\!\pmod{12}}
 \right)^2.
\]
El coeficiente relativo de la rama orientada es

\[
 \boxed{\eta_{30r}^{\Omega}(\Gamma)
 =A_\Gamma(r)-A_{\Gamma_e}(r).}
\]
Entonces

\[
 4\eta_{90}^{\Omega}-3\eta_{120}^{\Omega}
 =\Omega(\tau_\Gamma)-80,
\]
por lo que esta capa prolonga el primer componente de \(K_\Omega\). Las dos familias forman el levantamiento bidireccional

\[
 \mathcal F_{\rm tower}^{2w}:
 \mathcal R_{324}\longrightarrow
 \mathcal E_{\mathfrak S}\times\mathcal E_{\mathfrak S},
 \qquad
 \Gamma\longmapsto
 \bigl(\eta^D(\Gamma),\eta^\Omega(\Gamma)\bigr).
\]
\Needspace{7\baselineskip}
\subsection{Determinación finita de la torre infinita}
La periodicidad de la palabra temporal da

\[
 D_\Gamma(9(r+12))=D_\Gamma(9r).
\]
Si

\[
 T_\Gamma=\sum_{j=0}^{11}\tau_{\Gamma,j},
 \qquad r=12q+s,\quad 0\le s<12,
\]
entonces

\[
 A_\Gamma(12q+s)
 =A_\Gamma(s)+(12q^2+2qs)T_\Gamma^2.
\]
Por tanto, toda la torre queda determinada por doce valores de \(D_\Gamma\), doce valores de \(A_\Gamma\) y el entero \(T_\Gamma^2\). La infinitud de alturas no exige una infinitud de datos independientes.

La misma fórmula resuelve el diamante genealógico de altura 360. Si

\[
 h(a,b)=90a+120b,
\]
entonces

\[
 (a+4,b)\sim(a,b+3)
\]
porque \(3(a+4)+4b=3a+4(b+3)\). La genealogía \((a,b)\) se conserva hasta aplicar memoria y acarreo; sólo después desciende a la altura común.

\Needspace{7\baselineskip}
\subsection{Convergencia absoluta del carácter refinado sobre las torres \(90/120\)}
Se tiene

\[
 |\eta_h^D|\le108,\qquad
 |\eta_{30r}^{\Omega}|\le24r^2.
\]
Como \(0<q_+<q_-<1\),

\[
 |s_{30r}|\le2q_-^{30r}.
\]
De aquí se deduce

\[
 \sum_{h\in\mathfrak S}|\eta_h^D s_h|<\infty,
 \qquad
 \sum_{h\in\mathfrak S}|\eta_h^\Omega s_h|<\infty.
\]
El carácter refinado está bien definido sin truncar la torre.

\Needspace{7\baselineskip}
\subsection{Memoria, acarreo y cociclo}
El coeficiente desnudo se enriquece mediante

\[
 \mathbb F_{\rm tower}(\Gamma)_h
 =\bigl(\eta_h^D,\eta_h^\Omega,\rho_h,
 \mathbf w_h,c_h,N_h^{\rm surv}\bigr),
\]
donde la firma truncada satisface

\[
 \sigma_{\Gamma|\ell_h}
 =r(\rho_h)+\mathcal D\mathbf w_h,
 \qquad
 \mathcal D=\operatorname{diag}(2,6,2,2,4),
\]
y el acarreo de dos residuos es

\[
 c(\rho_1,\rho_2)=
 \mathcal D^{-1}
 \bigl[r(\rho_1)+r(\rho_2)-r(\rho_1+\rho_2)\bigr].
\]
El defecto de localidad de una composición,

\[
 \delta_h(\Gamma,\Lambda)=
 F_h(\Gamma\star\Lambda)-F_h(\Gamma)-F_h(\Lambda),
\]
obedece la identidad de cociclo

\[
 \delta_h(\Gamma,\Lambda)
 +\delta_h(\Gamma\star\Lambda,\Xi)
 =\delta_h(\Lambda,\Xi)
 +\delta_h(\Gamma,\Lambda\star\Xi).
\]
Así se preserva por qué dos rutas con la misma altura visible pueden transportar historias distintas.

\Needspace{7\baselineskip}
\subsection{Selección de la lectura por el acoplamiento}
Las dos ramas no se promedian por conveniencia. El acoplamiento físico \(a\) debe determinar una transformación natural

\[
 \Pi_a:
 \mathcal E_{\mathfrak S}\times\mathcal E_{\mathfrak S}
 \longrightarrow\mathcal E_{\mathfrak S}
\]
que conmute con truncaciones y simetrías, respete el cociclo de acarreo, conserve el origen electrónico y sea compatible con \(X_{90}^{4}=Y_{120}^{3}\). Si el acoplamiento intercambia las dos lecturas, la única clausura continua, multiplicativa, simétrica y normalizada sobre escalares positivos es la media geométrica; en exponentes,

\[
 L_{\rm sym}=\frac{L_D+L_\Omega}{2}.
\]
En ausencia de esa simetría, la salida es la medida espectral conjunta de \((L_D,L_\Omega)\), no una media arbitraria.

\Needspace{7\baselineskip}
\subsection{Refinamiento coinductivo por órbitas primitivas}
El levantamiento finito anterior admite una prolongación. Sea \(\mathscr X_\Gamma\) el grafo de extensiones admisibles de la ruta, \(U_\Gamma\) su operador de transferencia y \(R_\Gamma\) la involución de orientación. Las trazas

\[
 a_n(\Gamma)=\operatorname{Tr}(R_\Gamma U_\Gamma^n)
\]
y la inversión de Möbius

\[
 p_n(\Gamma)=\frac1n\sum_{d\mid n}\mu(d)a_{n/d}(\Gamma)
\]
separan órbitas primitivas. La diferencia

\[
 p_h(\Gamma)-p_h(\Gamma_e)
\]
es un refinamiento ulterior del coeficiente de altura \(h\), no una sustitución de las dos lecturas finitas.

\Needspace{7\baselineskip}
\subsection{Formulación histórica mediante autómata de extensión}
La firma \(\mathbb Z^5\) no agota la historia de una ruta. Para cada \(\Gamma\in\mathcal R_{324}\) se conservan:

\begin{itemize}
\item las 108 transiciones del registro genealógico;
\item los veinte bloques y sus prefijos;
\item las palabras \(w_6\);
\item las nueve comparaciones \(K\leftrightarrow K+9\);
\item la hoja, el acarreo y el devanado;
\item los primeros retornos de clase, órbita, celda y estado cinemático;
\item la ambigüedad acumulada y las tríadas decimales estabilizadas;
\item los perfiles \(K_D\) y \(K_\Omega\).
\end{itemize}
Estos datos permiten construir un transductor de torres sin abrir el residual de una masa.

\Needspace{7\baselineskip}
\subsection{Autómata local de extensión}
Sea \(\mathscr X_\Gamma\) el conjunto finito de estados locales compatibles con la ruta. Un estado conserva

\[
 x=(t,g,B,p,\varepsilon,c,w),
\]
donde \(t\) es el tiempo discreto, \(g\) la fase nonádica, \(B\) el bloque, \(p\) el prefijo, \(\varepsilon\) la hoja, \(c\) el acarreo y \(w\) el devanado. Existe una arista \(x\to y\) cuando \(y\) es una extensión admisible de \(x\), respeta el registro genealógico y supera la poda \(G_9\). La tabla de 36.864 composiciones de acarreo proporciona la ley local de suma; las tablas de retorno fijan la memoria.

Sea \(U_\Gamma\) el operador de transferencia normalizado de este grafo y \(R_\Gamma\) la involución de orientación. Se exige

\[
 U_{C_M\Gamma}=J U_\Gamma J^{-1},\qquad
 R_{C_M\Gamma}=-J R_\Gamma J^{-1}.
\]
\Needspace{7\baselineskip}
\subsection{Conteos orientados de celdas primitivas y construcción de la firma electrónica}
Para \(n\ge1\), la traza orientada

\[
 a_n(\Gamma)=\operatorname{Tr}(R_\Gamma U_\Gamma^n)
\]
cuenta retornos cerrados con signo. La inversión de Möbius separa órbitas primitivas:

\[
 p_n(\Gamma)
 =\frac1n\sum_{d\mid n}\mu(d)a_{n/d}(\Gamma).
\]
Se toma el electrón como origen afín de las correcciones relativas y se define

\[
 \boxed{
 \eta_h(\Gamma)=p_h(\Gamma)-p_h(\Gamma_e),
 \qquad h\in\mathfrak S.}
\]
Así se obtiene el mapa

\[
 \mathcal F_{\rm tower}:
 \Gamma^{\rm enr}\longmapsto
 \eta(\Gamma)=(\eta_h(\Gamma))_{h\in\mathfrak S}.
\]
No hay parámetros por especie. El mismo autómata, la misma involución y la misma inversión de Möbius actúan sobre las 324 rutas.

\Needspace{7\baselineskip}
\subsection{Convergencia del refinamiento primitivo}
Sea \(N_\Gamma=\dim\mathscr X_\Gamma\). Si \(U_\Gamma\) es una contracción,

\[
 |a_n(\Gamma)|\le N_\Gamma.
\]
El número de divisores de \(n\) crece sublinealmente, por lo que

\[
 |p_n(\Gamma)|
 \le\frac{N_\Gamma\tau(n)}n.
\]
Como \(0<q_+<q_-<1\),

\[
 |s_h|\le q_-^h+q_+^h\le2q_-^h.
\]
Por comparación,

\[
 \sum_{h\in\mathfrak S}|\eta_h(\Gamma)s_h|<\infty
\]
uniformemente en cualquier colección finita de fibras. El carácter refinado está, por tanto, bien definido.

\Needspace{7\baselineskip}
\subsection{Identidades de conservación de la ruta electrónica y unicidad del punto fijo central}
La construcción debe satisfacer simultáneamente:

\begin{enumerate}
\item invariancia por renumeración de estados;
\item covarianza bajo conjugación de hoja;
\item compatibilidad con suma directa de componentes;
\item igualdad al recomputar desde registro genealógico y desde el grafo certificado;
\item conservación del origen \(\eta(\Gamma_e)=0\);
\item insensibilidad absoluta a una permutación de las columnas externas de masa;
\item determinación de las colas mediante una cota anterior al contraste.
\end{enumerate}
La representación exacta de un residual dado demuestra completitud del codominio, pero no sustituye estas siete pruebas.

\Needspace{7\baselineskip}
\section{Reloj de ruta y normalización absoluta}
\Needspace{7\baselineskip}
\subsection{Los retornos enteros no bastan para distinguir especies}
En el atlas de 324 rutas, el primer retorno de celda es 27, el retorno cinemático dentro de CT--108 es 54 y el retorno cinemático extendido es 216. Al ser comunes, estos enteros fijan la arquitectura temporal global, pero no pueden explicar por sí solos diferencias de masa entre especies.

\Needspace{7\baselineskip}
\subsection{Holonomía temporal de la ruta electrónica y retorno con memoria}
La información específica reside en la holonomía del operador de extensión. Sea \(\lambda_j(\Gamma)\) un autovalor no trivial de \(U_\Gamma\) sobre el subespacio persistente. Escribimos

\[
 \lambda_j(\Gamma)=r_j(\Gamma)e^{i\theta_j(\Gamma)}.
\]
La frecuencia interna asociada es

\[
 \omega_j(\Gamma)
 =\frac{\theta_j(\Gamma)+2\pi w_j(\Gamma)}{108t_0},
\]
donde el entero \(w_j\) lo fija el devanado registrado, no una elección de rama posterior. Para un modo estable \(r_j=1\); para una resonancia \(r_j<1\) y su atenuación aporta una anchura.

El cociente

\[
 \rho_{\omega,j}(\Gamma)
 =\frac{\omega_j(\Gamma)}{\omega_0}
\]
es adimensional y puede modificar razones de masa sin alterar la década común \(-34\). La energía y la masa del modo son

\[
 E_j(\Gamma)=\hbar_{\rm ret}\omega_j(\Gamma)
 \chi_{\rm full}(\Gamma)R_{\rm act}^{\kappa_j(\Gamma)},
\]
\[
 m_j(\Gamma)=E_j(\Gamma)c_{\rm int}^{-2}.
\]
Esta ecuación reúne acción, reloj, carácter y lector sin convertir ninguno de ellos en otro. Su aplicación numérica exige determinar \(U_\Gamma\), fijar el devanado y declarar la carta dimensional.

\Needspace{7\baselineskip}


\endgroup

\begingroup
\let\chapter\section
\let\section\subsection
\let\subsection\subsubsection
\let\subsubsection\paragraph
\let\clearpage\relax
% Vista científica exacta materializada desde una rama condicional.
% Fuente: source/demostraciones_primarias/14_05b_extension_superviviente_caracter.tex; rama: HMTIncludeRouteLedgerBlock.
% No modifica el contenido matemático de la rama de origen.
\chapter{Extensión superviviente, registro dodecafásico y firma integral}
\label{chap:extension-ruta-caracter}

La trayectoria que se observa no constituye el objeto primario de esta
construcción. El objeto primario es una extensión compatible a todas las
profundidades del transporte HMT; una ruta finita es una de sus lecturas. Esta
distinción permite formular sin ambigüedad la cadena interna que conduce desde
una historia enriquecida hasta el carácter de acción. También determina qué
información puede olvidarse sin alterar la salida y cuál debe permanecer en el
registro.

El propósito de la sección es reunir en un solo argumento cuatro piezas que se
habían estudiado por separado: la extensión superviviente, la memoria
intrínseca, el registro orientado y la firma entera. El resultado terminal de
esta sección es el carácter formal
\[
 \Gamma_{108}^{\rm enr}
 \xrightarrow{\ L_{\rm tick}\ }
 \ZZ^5
 \xrightarrow{\ \chi_{\rm form}\ }
 \operatorname{Fun}(\mathcal P,\RR_{>0}),
\]
anterior tanto a la especialización de los coeficientes como a toda unidad
física. El primer mapa extrae la firma del registro leído; el segundo la
convierte en una función positiva sobre el espacio de parámetros
\(\mathcal P\). La especialización numérica sólo se efectúa después de
construir los coeficientes globales en la sección correspondiente. La
realización metrológica pertenece a una carta dimensional todavía posterior.
El mapa \(L_{\rm tick}\) es el lector integral de iteración y retorno del
registro cronológico enriquecido.

\section{Composición operatoria de la prolongación y cocadena torsional en la generación del registro de ruta}

La ley de masas comienza antes de la fórmula exponencial. APP aporta un único
soporte de rutas y sus dos hojas, suma y producto. El TRIT orienta cada
iteración elemental.
Si
\[
 r_t=1+((t-1)\bmod9),\qquad m_t:=M_{\rm ph}(r_t),
\]
el selector de fase \(M_{\rm ph}\) activa
\[
 \operatorname{op}_t=
 \begin{cases}
  \Sigma,&m_t=+1,\\
  \Pi,&m_t=-1,\\
  \varnothing\quad(\text{\rm retención o contador neutral}),&m_t=0.
 \end{cases}
\]
La lectura temporal puede interpretar el tercer caso como umbral o presente,
pero no lo convierte en una tercera hoja aritmética. No es APP quien alterna
sus hojas. TPK transporta
conjuntamente posición, orientación, hoja, fase, bloque, acarreo, prefijo,
cilindro, genealogía y memoria. El registro cronológico enriquecido anota
cada transición ---posición,
hoja, lectura, acarreo, fase, bloque y evento---. La firma local de frontera
se deriva después mediante \(\operatorname{Sig}_q(h,c,\lambda)\), y la firma
global de masa se extrae del registro completo; ninguna se reconstruye
retrospectivamente a partir de una masa.

El embudo finito que alimenta esta historia conserva cuatro tipos distintos:
\[
 9^2\!\cdot9^2\!\cdot4\!\cdot4=104\,976
 \longrightarrow 468\ \text{emisiones }U_6
 \longrightarrow 243\ \text{palabras }w_6
 \hookrightarrow \mathbb F_3^6,
 \qquad |\mathbb F_3^6|=729.
\]
La primera cifra cuenta condiciones iniciales del motor de dos cursores; 468
cuenta emisiones decimales y 243 palabras ternarias visibles. Las 324 semillas
de la proyección cronológica \(\Pi_{108}^{\rm cal}\) de un solo cursor
tampoco sustituyen las 104\,976
condiciones del estado completo. Cada flecha cambia, por tanto, de dominio y de
operación. La igualdad numérica entre el ambiente de \(729\) palabras y los
\(729\) sufijos de una etapa de filtración sólo se relaciona mediante la codificación
ternaria declarada; no identifica ambos objetos.

La elevación conserva asimismo su orden causal:
\[
 w_6\xrightarrow{\,L_0\,}w_{12}
 \xrightarrow{\,L_0\,}w_{18}
 \xrightarrow{\,L_1\,}w_{24}
 \xrightarrow{\,L_1\,}w_{30}
 \longrightarrow R_{36}\longrightarrow\mathcal G_9
 \longrightarrow
 \text{sección superviviente}.
\]
La frontera \(R_{36}\) abre la prolongación. En cada etapa de filtración,
\(\mathcal G_9\) agota
los 729 subcilindros y conserva el único compatible con todos los horizontes; el
episodio especular profundo \(3\to1\) es otra poda y no debe confundirse con
esa partición ordinaria \(729\to1\). El retorno nonádico conserva la fase,
pero cambia bloque, cilindro y memoria. Ni cifras objetivo ni masas intervienen
en esta selección. En la frontera dimensional, una partícula es una sección
persistente enriquecida de la extensión; la ruta es su historia observable.
Su nombre físico se lee después y no decide qué rama sobrevive.

Antes de abrir el registro cronológico dimensional, la misma genealogía conserva una
ventana publicada de veinte sextetos por canal y admite la primera lectura
arquimediana de doce tríadas decimales. Estos dos testigos pertenecen a la
prolongación del estado: no son los doce eventos del registro de masa, ni los
reemplazan. Los eventos se construyen sólo después, al leer la ruta sobre el
sistema dodecafásico orientado.

La genealogía dodecafásica convierte después los doce eventos orientados en un
registro y sólo entonces en la firma
\[
 \gamma_x\longmapsto\mathcal L(\gamma_x),\qquad
 \sigma_{\rm reg}:=L_{\rm tick}(\gamma_x)
 =(n_A,s_C,k_{\Dfour},b_{120},b_{\zeta}),
\]
\[
 \sigma_{\rm reg}\longmapsto
 \chi_{\rm form}(\sigma_{\rm reg};\,\cdot).
\]
Las coordenadas públicas \(b_{120}:=\nu_{120}^{\rm ev}\) y
\(b_{\zeta}:=\nu_{270}\) son salidas del registro cronológico, no correcciones
metrológicas introducidas al final. Los coeficientes posteriores del
semigrupo de torres conservan otra procedencia:
\[
 N_{120}=b_{120}+\eta_{120},\qquad
 b_{\zeta}\zeta_{270}\not\equiv\eta_{270}s_{270},
\]
con \(\zeta_{270}=135\Dfour^2\) y
\(s_{270}=q_-^{270}-q_+^{270}\). El refinamiento de torres se añade después
del carácter central sin borrar la procedencia de ninguna coordenada.

Tres salidas laterales conservan también su tipo. El sector nulo abandona la
rama positiva mediante su lector y su carta nulos; no es el límite
\(m\to0\) de un carácter positivo. Una trayectoria virtual es una extensión
finita con obstrucción terminal y conserva registro cronológico y acarreo sin adquirir por
ello una masa imaginaria. Finalmente, la sombra Witt acompaña cada etapa de filtración y
preserva la genealogía excepcional, pero no selecciona la ruta ni los
coeficientes de masa.

\begin{figure}[p]
\centering
\begin{tikzpicture}[
  font=\small,
  node distance=3.5mm,
  box/.style={draw,rounded corners=1.5mm,align=center,text width=12.8cm,
    inner xsep=6pt,inner ysep=4pt},
  hmt/.style={box,draw=hmtblue,fill=hmtlightblue},
  md/.style={box,draw=hmtviolet,fill=hmtlightviolet},
  act/.style={box,draw=hmtgreen,fill=hmtlightgreen},
  result/.style={box,draw=hmtgold,fill=hmtlightgold},
  arr/.style={-{Stealth[length=2.2mm]},line width=.7pt}]

  \node[hmt] (app) {\textbf{APP}: toro aritmético nonádico, caminos y hojas
    aditiva y multiplicativa. La especificación completa contiene
    \(104\,976\) semillas de doble cursor.};
  \node[hmt,below=of app] (trit) {\textbf{TRIT}: orientación local de cada
    transición; el levantamiento conserva residuo, cociente y acarreo.};
  \node[hmt,below=of trit] (tpk) {\textbf{TPK}: transporte de fase con memoria;
    \(104\,976\to468\,U_6\to243\,w_6\subset\mathbb F_3^6\).};
  \node[hmt,below=of tpk] (surv) {Filtración coinductiva y monodromía
    \(\mathcal G_9\): trayectoria superviviente con genealogía conservada.};
  \node[md,below=of surv] (route) {La trayectoria de partícula produce,
    sin consultar la masa, un registro dodecafásico y la firma
    \(\sigma=(n_A,s_C,k_{\Delta_4},b_{120},b_\zeta)\in\mathbb Z^5\).};
  \node[md,below=of route] (chi) {El carácter central
    \(\chi_0(\sigma)\) y el refinamiento armónico
    \(\eta\in\mathbb Z^{\langle90,120\rangle}\) construyen
    \(\chi_{\rm full}(\sigma,\eta)\in\mathbb R_{>0}\).};
  \node[act,below=of chi] (action) {Rama dimensional independiente:
    \(\Sigma_H=\varphi\alpha_{\rm HMT}^{16}\to\hbar_{\rm int}\),
    \(\omega_0=2\pi/(108t_0)\),
    \(m_0=\hbar_{\rm int}\omega_0c_{\rm int}^{-2}\).};
  \node[result,below=of action] (mass) {Confluencia tipada:
    \(m_X^{\rm int}=m_0\,\chi_{\rm full}(\sigma_X,\eta_X)\).
    La escala común fija la normalización absoluta y se cancela en
    \(m_X/m_Y\).};
  \node[result,below=of mass] (compare) {La carta de unidades publica la magnitud.
    CODATA y PDG se abren únicamente en la columna final de contraste.};

  \foreach \a/\b in {app/trit,trit/tpk,tpk/surv,surv/route,route/chi,
    chi/action,action/mass,mass/compare}{\draw[arr] (\a)--(\b);}
\end{tikzpicture}

\caption{Derivación de una magnitud de masa a partir de una trayectoria HMT.
El carácter aporta una razón adimensional; la escala de acción y el reloj
aportan la sección dimensional. La referencia externa interviene sólo después
de aplicar la carta de unidades.}
\label{fig:derivacion-masa-rev9}
\end{figure}

\section{Proyección de una extensión superviviente sobre su ruta observable}

Sea \((\mathcal S_m,\rho_{m+1,m})_{m\geq0}\) el sistema de estados HMT
supervivientes a profundidad \(m\), con mapas de truncamiento
\(\rho_{m+1,m}:\mathcal S_{m+1}\to\mathcal S_m\). Definimos
\[
 \Omega_{\HMT}^{\rm surv}
 :=\varprojlim_m\mathcal S_m
 =\left\{(x_m)_m:\rho_{m+1,m}(x_{m+1})=x_m\right\}.
\]
Un elemento \(x\in\Omega_{\HMT}^{\rm surv}\) conserva simultáneamente la
compatibilidad de todos sus truncamientos. Una carta de lectura de profundidad
108 produce
\[
 \operatorname{sh}_{108}(x)=\gamma_x\in\Gamma_{108}^{\rm enr}.
\]
La notación \(\operatorname{sh}\) recuerda que \(\gamma_x\) es una sombra
observable: registra una presentación secuencial del estado, pero el estado no
se define escogiendo retrospectivamente esa ruta.

La presentación enriquecida conserva, además de la palabra visible, fase,
hoja, orientación, borde, memoria antihoja y los predecesores necesarios para
evaluar los cociclos. Dos palabras con los mismos cinco enteros finales pueden
ser presentaciones distintas porque esos enteros se calculan después de
aplicar signos y correlaciones.

\begin{definition}[Memoria intrínseca y registro leído]
Una extensión dodecafásica enriquecida es una clase compatible de
truncamientos. Denotamos por \(\mathscr M(x)\) la memoria intrínseca que esos
truncamientos conservan. Fijada la carta
\(\gamma_x=\operatorname{sh}_{108}(x)\), su registro leído es
\[
 \mathcal L(\gamma_x)
 =\bigl(n_A,e_0,\ldots,e_{11},T_\zeta,C_\zeta\bigr),
\]
donde \(n_A\) es el conteo del soporte de acción declarado, los \(e_m\) son
los doce eventos orientados y \(T_\zeta,C_\zeta\) son la memoria torsional
total y su restricción a la corona. La compatibilidad exige que estos datos
sean preservados por todo truncamiento que contenga los pasos empleados en su
definición.
\end{definition}

La extensión ya contiene la memoria que la ruta presenta; la ruta no la crea.
Relativamente a la carta declarada, la composición
\[
x\longmapsto\gamma_x
\longmapsto\mathcal L(\gamma_x)
\longmapsto L_{\rm tick}(\gamma_x)
\]
es una aplicación. La invariancia frente a otra carta de presentación
requeriría, además, demostrar la compatibilidad de los correspondientes mapas
de transición. En este sentido preciso, el registro no se añade a una
partícula ya construida: es una lectura de la memoria de la extensión.

\section{Descomposición integral de \(12\) eventos en los sectores \(1+5+6\)}

Sea \(d_1,\ldots,d_{108}\in\{1,\ldots,9\}\) la palabra de una presentación
enriquecida. En la convención de paridad usada por el extractor, el soporte de
acción es
\[
 \mathcal D_{\rm act}=\{1,3,5,7\},
 \qquad
 n_A=\#\{t:d_t\in\mathcal D_{\rm act}\}.
\]
La elección de este soporte es parte declarada de la carta; no debe
confundirse con la clasificación TRIT de los pasos. Para los eventos pares se
usa
\[
 \mathcal D_{\rm evt}=\{2,4,6,8\},
 \qquad
 \Sigma_{12}=(+,+,+,-,-,-,+,+,+,-,-,-).
\]
El evento orientado del bloque \(m\in\{0,\ldots,11\}\) es
\[
 e_m=\Sigma_{12}(m)
 \#\{t\in\{9m+1,\ldots,9m+9\}:d_t\in\mathcal D_{\rm evt}\}.
\]

Emparejemos las dos semicoronas mediante
\[
 p_i=e_i+e_{i+6},
 \qquad
 a_i=e_i-e_{i+6},
 \qquad 0\leq i<6.
\]
Definimos
\[
 s_C=\sum_{i=0}^{5}p_i,
 \qquad
 M_i=p_i-p_5\quad(0\leq i<5),
\]
y escribimos \(M_5=(M_0,\ldots,M_4)\),
\(A_6=(a_0,\ldots,a_5)\).

El \cref{thm:memoria-1-5-6}, en la residencia primaria del sistema
dodecafásico de TPK, demuestra que
\[
 (e_0,\ldots,e_{11})\longmapsto(s_C,M_5,A_6)
\]
es inyectivo sobre su imagen y proporciona la inversa
\[
 p_5=\frac{s_C-\sum_{i=0}^{4}M_i}{6},\qquad
 p_i=p_5+M_i,
\]
\[
 e_i=\frac{p_i+a_i}{2},\qquad
 e_{i+6}=\frac{p_i-a_i}{2}.
\]
Aquí sólo se recuperan esas coordenadas para construir el carácter; la
prueba y la identidad \(1+5+6=12\) no se duplican.

El dato \(s_C\) es el que entra en el carácter angular. Las once coordenadas
\((M_5,A_6)\) conservan la colocación de ese total. Eliminarlas antes de
componer historias confunde eventos diferentes que pueden producir
correlaciones distintas.

\section{Cocadena torsional a nivel de paso}

El conteo de eventos no determina por sí solo la contribución torsional. Para
cada aparición de nueve, conservamos la cifra que la precede y definimos
\[
 j_2(d)=
 \begin{cases}
 0,&d=9,\\
 +1,&d\in\{2,4,6,8\},\\
 -1,&d\in\{1,3,5,7\}.
 \end{cases}
\]
Con índices cíclicos en la palabra, sea
\[
 \zeta_t=\mathbf 1_{\{d_t=9\}}j_2(d_{t-1}),
\]
\[
 T_\zeta=\sum_{t=1}^{108}\zeta_t,
 \qquad
 C_\zeta=\sum_{t\in\{27,54,81,108\}}\zeta_t.
\]
La coordenada fuerte que acompaña a \(\Dfour\) es
\[
 \boxed{k_{\Dfour}=T_\zeta-2C_\zeta.}
\]
La fórmula separa el arrastre acumulado de su corrección orientada en los
cuatro cierres de corona. El censo exhaustivo muestra que ni
\(T_\zeta\) ni \(C_\zeta\) por separado separan todas las fibras del censo;
el par ordenado sí lo hace en las dos convenciones tipadas que se contrastan.

Las dos coordenadas de torre restantes se calculan a partir del registro de
eventos. Sea
\[
 \tau_m=\operatorname{sgn}(e_m).
\]
Para las cuatro clases de paso cuatro,
\[
 I_a=\{a,a+4,a+8\}\pmod{12},
 \qquad 0\leq a<4,
\]
definimos el carácter ponderado por eventos
\[
 \nu_{120}^{\rm ev}
 =\sum_{a=0}^{3}\operatorname{sgn}
 \left(\sum_{m\in I_a}e_m\right),
\]
y la autocorrelación de paso tres
\[
 \nu_{270}=\sum_{m=0}^{11}\tau_m\tau_{(m+3)\bmod12}.
\]
La variante
\[
 \nu_{120}^{\rm sgn}
 =\sum_{a=0}^{3}\operatorname{sgn}
 \left(\sum_{m\in I_a}\tau_m\right)
\]
es otro observable. No se identifican silenciosamente: difieren en
25 de las 81 rutas del atlas.

\begin{definition}[Extractor enriquecido de firma]
Fijada la convención de eventos, se define
\[
 L_{\rm tick}(\gamma)=
 \bigl(n_A,s_C,k_{\Dfour},\nu_{120}^{\rm ev},\nu_{270}\bigr)
 \in\ZZ^5.
\]
El subíndice indica que la coordenada torsional se calcula antes de olvidar
los predecesores de los pasos nulos.
\end{definition}

Esta aplicación es determinista y no consulta masas, unidades ni etiquetas
de partículas. No se afirma que sea lineal respecto de una concatenación ya
proyectada: los signos y las autocorrelaciones obligan a componer primero los
registros enriquecidos.

\endgroup
\begingroup
\let\chapter\section
\let\section\subsection
\let\subsection\subsubsection
\let\subsubsection\paragraph
\let\clearpage\relax
% Vista científica exacta materializada desde una rama condicional.
% Fuente: source/demostraciones_primarias/14_05b_extension_superviviente_caracter.tex; rama: HMTIncludeCharacterBlock.
% No modifica el contenido matemático de la rama de origen.
\chapter{Carácter positivo de la firma y composición de registros}
\section{Carácter formal universal}

Sea \(\mathcal P=\RR^5\), con coordenadas
\[
 X=(X_A,X_C,X_4,X_{120},X_{270}).
\]
Para \(v=(n_A,s_C,k,\nu_{120},\nu_{270})\in\ZZ^5\), definimos
\[
 \ell_{\rm form}(v;X)
 =n_AX_A+s_CX_C+kX_4
  +\nu_{120}X_{120}+\nu_{270}X_{270},
\]
\[
 \chi_{\rm form}(v;X)=\exp\!\bigl(\ell_{\rm form}(v;X)\bigr).
\]
Así, \(\chi_{\rm form}\) toma valores en el grupo multiplicativo
\(\operatorname{Fun}(\mathcal P,\RR_{>0})\), con producto puntual. Todavía no
ha seleccionado coeficientes globales ni una especie.

\begin{theorem}[Carácter universal de la firma]
\label{thm:caracter-formal-ruta}
El mapa
\[
 \chi_{\rm form}:(\ZZ^5,+)\longrightarrow
 \bigl(\operatorname{Fun}(\mathcal P,\RR_{>0}),\cdot\bigr)
\]
es un homomorfismo. La composición
\[
 \chi_{\rm route}^{\rm form}
 :=\chi_{\rm form}\circ L_{\rm tick}
\]
evalúa formalmente cualquier registro enriquecido y no introduce una unidad
física.
\end{theorem}

\begin{proof}
Para todo \(X\in\mathcal P\), el exponente
\(\ell_{\rm form}(\,\cdot\,;X)\) es aditivo sobre \(\ZZ^5\). Por tanto,
\[
\chi_{\rm form}(v+w;X)
=\chi_{\rm form}(v;X)\chi_{\rm form}(w;X),
\]
y la identidad vale como igualdad de funciones positivas.
\end{proof}

El carácter formal precede a la construcción de los coeficientes globales.
En esta sección no se define todavía un carácter numérico de masa ni una
carta de unidades.

\section{Sección interna mínima sobre las 81 celdas}

El atlas finito contiene 81 celdas orientadas. Al proyectarlas sobre la carta
de familia de ruta se obtienen 56 valores. Denotemos esa proyección por
\[
 \mathcal R:\mathcal C_9\longrightarrow\mathcal R_{56}.
\]
El histograma exacto de sus fibras es
\[
 41\times1,
 \qquad10\times2,
 \qquad2\times3,
 \qquad1\times4,
 \qquad2\times5,
\]
y su suma es 81.

Fijados el origen y el orden lexicográfico de la carta, sea
\(\kappa(x)\in\{0,1,2,3,4\}\) el rango de \(x\) dentro de su fibra.

\begin{theorem}[Selector interno mínimo]
\label{thm:selector-interno-minimo}
La aplicación
\[
 x\longmapsto\bigl(\mathcal R(x),\kappa(x)\bigr)
\]
es inyectiva sobre las 81 celdas. Ningún alfabeto de memoria con menos de
cinco símbolos puede levantar \(\mathcal R\) inyectivamente sobre todo el
atlas.
\end{theorem}

\begin{proof}
Dentro de cada fibra, \(\kappa\) asigna rangos distintos; entre fibras, la
primera coordenada ya es distinta. Existe una fibra de tamaño cinco, así que
el principio del palomar obliga a toda elevación inyectiva a disponer de al menos
cinco valores de memoria. El rango \(0,\ldots,4\) alcanza esa cota.
\end{proof}

La canonicidad es relativa al origen y al orden declarados. \(\kappa\) no es
una carga, un sabor ni una masa. Las trece clases gruesas del atlas producen
multisecciones finitas; sólo la clase central posee el punto fijo \((5,5)\)
bajo la inversión central. Las otras doce requieren información orientada
para seleccionar un punto y no admiten un selector puntual equivariante desde
la clase gruesa sola.

\subsection{Inversión de Möbius y caracteres de ida y retorno}

Sobre la palabra de signos dodecafásica definimos
\[
 C_M(\tau)_m=-\tau_{11-m}.
\]
Esta operación combina la inversión del orden de lectura con el cambio de
hoja. Un carácter lineal
\(\ell_a(\tau)=\sum_ma_m\tau_m\), con
\(a_m=a_{11-m}\), satisface
\[
 \ell_a(C_M\tau)=-\ell_a(\tau),
\]
mientras que toda forma cuadrática simétrica compatible con la reversión
satisface
\[
 Q(C_M\tau)=Q(\tau).
\]
Así, las coordenadas lineales distinguen orientación de ida y retorno y las
cuadráticas conservan el contenido de correlación. Esta paridad explica por
qué \(s_C\) cambia de signo mientras \(\nu_{270}\) puede permanecer
invariante. En el atlas, las siete formas de palabra tienen multiplicidades
\[
 54,2,2,2,7,7,7;
\]
por ello la antihoja se realiza como correspondencia entre fibras y no como
una involución biyectiva obtenida ignorando sus multiplicidades.

\section{Composición orientada de registros antes de la proyección}

Sean \(\Gamma\) y \(\Lambda\) dos historias enriquecidas, \(B\) una
subhistoria común y \(g\) el transporte que identifica fase, hoja,
orientación y posición en la interfaz. En el nivel de registro se define
\[
 \boxed{
 \Gamma\star_B\Lambda
 =\bigl(n_\Gamma+n_\Lambda-n_B,
        e_\Gamma+g e_\Lambda-e_B\bigr).}
\]
La definición presupone también la identificación del borde, los
predecesores torsionales, la corona y la memoria antihoja.

\begin{proposition}[Obstrucción al descenso pentadimensional de la composición de registros]
\label{prop:no-descenso-empalme-cinco}
En general no existe una operación binaria sobre \(\ZZ^5\) que recupere
\(L_{\rm tick}(\Gamma\star_B\Lambda)\) a partir de
\(L_{\rm tick}(\Gamma)\) y \(L_{\rm tick}(\Lambda)\) solamente.
\end{proposition}

\begin{proof}
La operación compone los doce eventos antes de aplicar
\(\operatorname{sgn}\) y las autocorrelaciones. Dos historias pueden tener
la misma cinco-tupla y valores distintos de \(M_5\) o \(A_6\). Al componerlas
con una tercera historia, alguna suma
\(e_a+e_{a+4}+e_{a+8}\) puede cruzar cero en una y no en la otra, cambiando
\(\nu_{120}^{\rm ev}\). Los rombos finitos proporcionan testigos explícitos.
Por tanto, la proyección a cinco enteros no
es una congruencia para \(\star_B\).
\end{proof}

La proposición formaliza una regla operativa esencial: las extensiones se
componen; las firmas se leen después. La partícula compuesta tampoco se
obtiene sumando a ciegas dos filas de una tabla, sino componiendo sus memorias
compatibles y evaluando el resultado.

\section{Dependencia dodecafásica del carácter de persistencia}

Los motores reducidos de período 54 imponen
\[
 e_{m+6}=e_m.
\]
De aquí se siguen las congruencias
\[
 n_A\equiv0\pmod2,
 \qquad
 s_C\equiv0\pmod2,
 \qquad
 \nu_{270}\equiv0\pmod4.
\]
Las firmas de la reconstrucción escalar abreviada violan al menos una de
estas congruencias. Por tanto, su genealogía exige el registro enriquecido y
no procede de ese motor reducido con las reglas declaradas.

Este resultado no es una obstrucción a la dinámica TPK enriquecida. Demuestra
qué memoria necesita: una dodecafase genuina, con posible
\(A_6\ne0\), composición orientada y predecesores torsionales conservados. La
ablación evita atribuir a TPK una limitación que pertenece a una reducción
periódica concreta.

\section{Factorización por el cociente de registro y realización física}

En el alcance de esta sección se declara la relación
\[
\gamma\sim_{\mathcal L}\gamma'
\quad\Longleftrightarrow\quad
\text{\(\gamma,\gamma'\) presentan la misma extensión y }
\mathcal L(\gamma)=\mathcal L(\gamma').
\]
Denotamos por
\(\mathcal E_{108}^{\rm surv}:=
\Gamma_{108}^{\rm enr}/\!\sim_{\mathcal L}\)
el cociente relativo a esta carta. Sea
\[
 \pi_{\rm ext}:\Gamma_{108}^{\rm enr}\longrightarrow
 \mathcal E_{108}^{\rm surv}
\]
la proyección canónica. Por construcción, el extractor es constante en sus
fibras. La afirmación más fuerte según la cual toda carta admisible de una
misma extensión produce automáticamente el mismo registro requeriría demostrar
la compatibilidad entre cartas y no se presupone aquí.

\begin{theorem}[Factorización relativa por el cociente de registro]
\label{thm:descenso-caracter-extension}
Existe una única aplicación
\[
 \overline L:\mathcal E_{108}^{\rm surv}\longrightarrow\ZZ^5
\]
tal que
\[
 \boxed{L_{\rm tick}=\overline L\circ\pi_{\rm ext}.}
\]
Por tanto, \(\chi_{\rm form}\circ\overline L\) es una función de la clase en
el cociente relativo y no del representante de ruta escogido para leerla.
\end{theorem}

\begin{proof}
Si \(\pi_{\rm ext}(\gamma)=\pi_{\rm ext}(\gamma')\), ambas presentaciones
poseen el mismo registro por la relación declarada y, por las fórmulas
anteriores, la misma cinco-tupla. Definimos
\(\overline L([\gamma])=L_{\rm tick}(\gamma)\). La constancia prueba que la
definición es independiente del representante; la sobreyectividad del mapa
al cociente prueba la unicidad.
\end{proof}

La construcción interna llega así hasta
\[
 \Omega_{\HMT}^{\rm surv}
 \longrightarrow\mathcal E_{108}^{\rm surv}
 \xrightarrow{\ \overline L\ }\ZZ^5
 \xrightarrow{\ \chi_{\rm form}\ }
 \operatorname{Fun}(\mathcal P,\RR_{>0}).
\]
El contenido del teorema es positivo y completo en este dominio: una
extensión superviviente determina su clase de registro, la clase determina
una única firma y la firma determina el carácter formal. La especialización
global y la jerarquía completa de torres se construyen en los capítulos
siguientes; las
comparaciones con nombres y unidades físicas se presentan después, sin
alterar esta cadena.

\endgroup
\begingroup
\let\chapter\section
\let\section\subsection
\let\subsection\subsubsection
\let\subsubsection\paragraph
\let\clearpage\relax
\chapter{Álgebra partida y factorización exacta del carácter}
\label{chap:factorizacion-split}
\index{álgebra partida}\index{carácter!factorización partida}
\index{cono nulo}\index{canales $q_\pm$}

La firma formal de una extensión y los canales angulares proceden de ramas
ya construidas. En esta sección ambas confluyen en una factorización
determinantal que conserva, por separado, la escala radial y la orientación
relativa. El resultado pertenece al estrato algebraico: no identifica por sí
solo los ideales nulos con fotones físicos ni el interior positivo con una
partícula material.

\section{Coordenadas angulares conjugadas y retícula espectral bidireccional}
\label{sec:reticula-angular-split}

La cadena causal anterior fija
\[
 A_{\rm rad}=\frac{\pi A}{180},\qquad
 \Cstarrad=\frac{\pi\Cstar}{180},
\]
y, sólo después de obtener la coordenada directa \(A\) y la coordenada
conjugada única \(C^*\),
\[
 q_+=\exp[-(A_{\rm rad}+\Cstarrad)],
 \qquad
 q_-=\exp[-(A_{\rm rad}-\Cstarrad)].
\]
Para una firma angular \((n_A,s_C)\), introducimos los dos números de
enrollamiento
\[
 W_+=6n_A+s_C,
 \qquad
 W_-=6n_A-s_C.
\]

\begin{theorem}[Retícula angular de índice doce]
\label{thm:reticula-two-way}
La aplicación
\[
 \mathcal W:\ZZ^2\to\ZZ^2,
 \qquad
 \binom{n_A}{s_C}\longmapsto
 \binom{W_+}{W_-}
 =\begin{pmatrix}6&1\\6&-1\end{pmatrix}
 \binom{n_A}{s_C}
\]
es inyectiva y su imagen tiene índice doce. Sus divisores invariantes son
\((1,12)\), y
\[
 \operatorname{im}\mathcal W
 =\{(u,v)\in\ZZ^2:u+v\equiv0\pmod{12}\}.
\]
Además,
\[
 \boxed{
 \exp\!\left(n_AA_{\rm rad}+\frac{s_C}{6}\Cstarrad\right)
 =\left(q_+^{-W_+}q_-^{-W_-}\right)^{1/12}.}
\]
\end{theorem}

\begin{proof}
La matriz tiene determinante \(-12\). Como el máximo común divisor de sus
entradas es uno, su forma normal de Smith es
\(\operatorname{diag}(1,12)\). Si
\((u,v)=(6n+s,6n-s)\), entonces \(u+v=12n\). Recíprocamente, si
\(u+v\) es múltiplo de doce, entonces
\[
 n=\frac{u+v}{12},
 \qquad
 s=\frac{u-v}{2}
\]
son enteros y recuperan \((u,v)\). Finalmente,
\[
 \frac{W_+(A_{\rm rad}+\Cstarrad)
       +W_-(A_{\rm rad}-\Cstarrad)}{12}
 =n_AA_{\rm rad}+\frac{s_C}{6}\Cstarrad,
\]
lo que prueba la identidad al exponenciar.
\end{proof}

La involución \(s_C\mapsto-s_C\) intercambia
\(W_+\leftrightarrow W_-\) y \(q_+\leftrightarrow q_-\). La raíz
duodécima es el índice integral de la imagen de \(\mathcal W\), no un
redondeo. La coordenada \(C^*\) actúa como componente antisimétrica de la retícula;
no es una corrección decimal de una masa ni una salida de
\(\Gamma_{6\to8}\).

\section{Álgebra real partida}
\label{sec:algebra-real-partida}

\begin{definition}[Álgebra real partida]
Sea
\[
 \mathcal A_{\rm sp}
 :=\RR[\mathsf R]/(\mathsf R^2-1).
\]
La conjugación y la norma se definen, respectivamente, por
\[
 \overline{a+b\mathsf R}=a-b\mathsf R,
 \qquad
 N_{\rm sp}(z)=z\overline z=a^2-b^2.
\]
\end{definition}

Los idempotentes conjugados
\[
 P_\pm=\frac{1\pm\mathsf R}{2}
\]
satisfacen
\[
 P_\pm^2=P_\pm,\qquad
 P_+P_-=0,\qquad
 P_++P_-=1,\qquad
 \overline{P_+}=P_-.
\]

\begin{proposition}[Descomposición espectral partida]
\label{prop:descomposicion-espectral-partida}
Todo \(z\in\mathcal A_{\rm sp}\) admite una descomposición única
\[
 z=r_+P_++r_-P_-,
 \qquad r_\pm\in\RR,
\]
y su norma es
\[
 \boxed{N_{\rm sp}(z)=r_+r_-.}
\]
Cada ideal principal \(\mathcal A_{\rm sp}P_\pm=\RR P_\pm\) es una recta
nula; en particular,
\[
 N_{\rm sp}(\lambda P_\pm)=0.
\]
\end{proposition}

\begin{proof}
Si \(z=a+b\mathsf R\), la elección \(r_+=a+b\), \(r_-=a-b\) proporciona la
descomposición, cuya unicidad se sigue de la independencia de
\(P_+\) y \(P_-\). La conjugación intercambia los dos coeficientes:
\(\overline z=r_-P_++r_+P_-\). La ortogonalidad de los idempotentes da
\[
 z\overline z=(r_+r_-)(P_++P_-)=r_+r_-.
\]
La afirmación sobre cada ideal se obtiene anulando uno de los dos
coeficientes.
\end{proof}

\begin{definition}[Cono causal partido]
\label{def:cono-causal-partido}
El cono positivo, su frontera nula y su interior son
\[
 \mathcal C_{\rm sp}^{+}
 =\{r_+P_++r_-P_-:r_+,r_-\geq0\},
\]
\[
 \partial_{\rm null}\mathcal C_{\rm sp}^{+}
 =\RR_{\geq0}P_+\cup\RR_{\geq0}P_-,
\qquad
 \operatorname{int}\mathcal C_{\rm sp}^{+}
 =\{r_+P_++r_-P_-:r_+r_->0\}.
\]
\end{definition}

La frontera posee una sola dirección activa y norma cero. El interior exige
dos direcciones simultáneas y norma positiva. Esta estratificación es exacta;
su realización física exige mapas que preserven la norma, el transporte y la
dinámica pertinentes.

\section{Transportador relativo y condiciones de realización representacional}
\label{sec:transportador-relativo-split-md}

El bloque histórico \(B_0=7I+2\mathsf R\) tiene espectro \(\{9,5\}\), y el
bloque angular \(B_1=AI+\Cstar\mathsf R\) tiene espectro
\(\{A+\Cstar,A-\Cstar\}\). El
\cref{thm:no-go-entrelazador-bloques-split} demuestra que, con los valores internos
vigentes, todo supuesto entrelazador lineal
\(\Phi B_0=B_1\Phi\) es nulo. El objeto activo es el transportador
multiplicativo relativo ya construido, no una equivalencia entre
representaciones de espectros disjuntos.

\section{Factorización partida del carácter angular}
\label{sec:factorizacion-split-angular}

Sea
\[
 \sigma=\frac{s_C}{6},\qquad
 w_\pm=\frac{n_A\pm\sigma}{2}=\frac{W_\pm}{12},
 \qquad
 r_\pm=q_\pm^{-w_\pm},
\]
y definamos
\[
 Z_v=r_+P_++r_-P_-.
\]

\begin{theorem}[Factorización partida del carácter angular]
\label{thm:split-character}
Para toda firma angular \(v=(n_A,s_C)\in\ZZ^2\),
\[
 \boxed{
 N_{\rm sp}(Z_v)
 =\det Z_v
 =\exp\!\left(
 n_AA_{\rm rad}+\frac{s_C}{6}\Cstarrad
 \right).}
\]
\end{theorem}

\begin{proof}
Por la definición de los canales,
\[
 \log r_+
 =\frac{n_A+\sigma}{2}(A_{\rm rad}+\Cstarrad),
 \qquad
 \log r_-
 =\frac{n_A-\sigma}{2}(A_{\rm rad}-\Cstarrad).
\]
Al sumar,
\[
 \log(r_+r_-)
 =n_AA_{\rm rad}+\sigma\Cstarrad.
\]
La \cref{prop:descomposicion-espectral-partida} da
\(N_{\rm sp}(Z_v)=r_+r_-\), y la representación diagonal en la base
\((P_+,P_-)\) da la misma expresión para el determinante.
\end{proof}

La componente que no ve el determinante conserva la orientación relativa:
\[
 \frac12\log\frac{r_+}{r_-}
 =\frac12\left(
 n_A\Cstarrad+\frac{s_C}{6}A_{\rm rad}
 \right),
\]
\[
 Z_v=
 \exp\!\left[
 \frac12\left(n_AA_{\rm rad}+\frac{s_C}{6}\Cstarrad\right)
 \right]
 \exp\!\left[
 \frac12\left(n_A\Cstarrad+\frac{s_C}{6}A_{\rm rad}\right)
 \mathsf R
 \right].
\]
El primer factor fija la escala radial; el segundo registra la orientación
relativa. Son dos lecturas coordinadas de un mismo elemento.

\section{Incorporación central de las torres}
\label{sec:factorizacion-split-torres}

Una vez construidos globalmente \(\Dfour\), \(s_{120}\) y
\(\zCorr=\zeta_{270}^{\rm corr}\), sea
\[
 t(v)=
 k_{\Dfour}\Dfour
 +\nu_{120}s_{120}
 +\nu_{270}\zCorr,
 \qquad
 \widetilde Z_v=e^{t(v)/2}Z_v.
\]

\begin{corollary}[Carácter completo como norma partida]
\label{cor:split-character-torres}
Fijados los bloques globales anteriores,
\[
 \boxed{
 N_{\rm sp}(\widetilde Z_v)
 =\exp\!\left(
 n_AA_{\rm rad}+\frac{s_C}{6}\Cstarrad
 +k_{\Dfour}\Dfour+\nu_{120}s_{120}
 +\nu_{270}\zCorr
 \right).}
\]
\end{corollary}

\begin{proof}
En la representación bidimensional, la multiplicación por el escalar central
\(e^{t(v)/2}\) multiplica el determinante por \(e^{t(v)}\). Se aplica el
\cref{thm:split-character}.
\end{proof}

El corolario construye la realización partida declarada del carácter activo;
no clasifica todas las realizaciones posibles ni transforma las torres en
pruebas independientes. Sus hipótesis son los canales y los bloques globales
ya fijados, y su control negativo es inmediato: omitir el factor central
\(1/2\) duplica el exponente de torre; activar simultáneamente ambos
idempotentes con coeficientes positivos abandona la frontera nula.

\endgroup
\begingroup
\let\chapter\section
\let\section\subsection
\let\subsection\subsubsection
\let\subsubsection\paragraph
\let\clearpage\relax
\chapter{Valores, completitud y bloques de la jerarquía espectral}
\label{chap:torres}

El operador bicapa, sus momentos y el semigrupo
\(\mathfrak S=\langle90,120\rangle\) se definen una sola vez en el
\cref{sec:jerarquia-espectral-infinita-rev7}. Este capítulo publica valores,
prueba la completitud reconstructiva y separa los bloques globales. Los
mismos enteros \(90,120\) aparecen también como cardinalidades en la
incidencia hexada--octada, pero esa coincidencia no transporta la incidencia
\(\iota_{6,8}\) a la carta angular ni convierte sus configuraciones en
momentos de masa.

\section{Primeras evaluaciones de la jerarquía}

Los índices
\[
 90,120,180,210,240,270,360,420
\]
son una selección de primeras alturas de \(\mathfrak S\), pues
\[
 180=2\cdot90,\quad210=90+120,\quad240=2\cdot120,
\]
\[
 270=3\cdot90,\quad360=3\cdot120,
 \quad420=2(90+120).
\]
Los momentos \(s_n=q_-^n-q_+^n\) y la interfaz generadora
\(\mathcal K_{\mathrm{ang}}^{90,120}=(s_{90},s_{120})\) se usan aquí con la
convención ya fijada; no se redefinen. La totalidad de los valores queda
determinada por \(A\), \(\Cstar\) y la recurrencia cuadrática del lector.

\begin{table}[htbp]
\centering
\caption{Momentos orientados de la jerarquía global.}
\label{tab:torres}
\begin{tabular}{@{}rr@{}}
\toprule
\(n\)&\(s_n\)\\
\midrule
90  & \(2.9516943928982796\times10^{-4}\)\\
120 & \(1.9683511250294992\times10^{-5}\)\\
180 & \(8.7345871869508335\times10^{-8}\)\\
210 & \(5.8181313057291163\times10^{-9}\)\\
240 & \(3.8754674969305336\times10^{-10}\)\\
270 & \(2.5814553607509555\times10^{-11}\)\\
360 & \(7.6293257871406076\times10^{-15}\)\\
420 & \(3.3850663628348497\times10^{-17}\)\\
\bottomrule
\end{tabular}
\end{table}

\begin{proposition}[Carácter global de los momentos]
Las cantidades \(s_n\) de la \cref{tab:torres} son invariantes de la carta
angular fijada. Una especie aporta coordenadas enteras que ponderan esos
invariantes, pero no modifica \(q_\pm\) ni introduce una torre propia.
\end{proposition}
\begin{proof}
Cada \(s_n\) es una función de \(q_-\) y \(q_+\). El teorema de inversión de
canales demuestra que la aplicación conjunta
\((A,\Cstar)\mapsto(q_+,q_-)\) es inyectiva. No aparece ninguna etiqueta de
especie en estas definiciones.
\end{proof}

\begin{theorem}[Completitud reconstructiva de la jerarquía]
\label{thm:completitud-reconstructiva-torres}
Sea \(n_1<n_2<\cdots\) la enumeración creciente sin repeticiones de
\(\mathfrak S\). Para
\[
 a_k:=s_{n_k}=q_-^{n_k}-q_+^{n_k}>0,
\]
toda razón positiva \(\rho\in\RR_{>0}\) admite una expansión absolutamente
convergente en la jerarquía:
\[
 \boxed{\rho=
 \exp\!\left(\sum_{k\geq1}\nu_k a_k\right)}
 \qquad\text{para alguna sucesión }(\nu_k)_{k\geq1}\subset\ZZ.
\]
\end{theorem}

\begin{proof}
Como \(0<q_+<q_-<1\), se cumple
\[
 0<a_k<q_-^{n_k}.
\]
Los \(n_k\) son múltiplos de \(30\) no menores que \(90\); por tanto,
\[
 \sum_{k\geq1}a_k
 \leq\sum_{m\geq3}q_-^{30m}
 =\frac{q_-^{90}}{1-q_-^{30}}<\infty,
 \qquad a_k\longrightarrow0.
\]
Fijemos \(r_0=\log\rho\) y una regla de desempate para el entero más próximo.
Definimos recursivamente
\[
 \nu_k=\operatorname{nint}\!\left(\frac{r_{k-1}}{a_k}\right),
 \qquad
 r_k=r_{k-1}-\nu_k a_k.
\]
La propiedad del entero más próximo da
\[
 |r_k|\leq\frac{a_k}{2}\longrightarrow0.
\]
La identidad telescópica
\[
 r_0=\sum_{j=1}^{N}\nu_j a_j+r_N
\]
implica \(r_0=\sum_{j\geq1}\nu_j a_j\). La convergencia es absoluta, pues
\[
 \begin{aligned}
 \sum_{k\geq1}|\nu_k|a_k
 &=\sum_{k\geq1}|r_{k-1}-r_k|\\
 &\leq |r_0|+\frac{a_1}{2}
 +\frac12\sum_{k\geq2}(a_{k-1}+a_k)<\infty.
 \end{aligned}
\]
Al exponenciar se obtiene la afirmación.
\end{proof}

\begin{remark}[Alcance y control negativo]
El teorema es una sobreyectividad reconstructiva del lenguaje de torres sobre
\(\RR_{>0}\): prueba resolución del codominio, no selección física. El
algoritmo recibe \(r_0=\log\rho\); si \(\rho\) es una magnitud objetivo, sus
coeficientes son, por ello, una reconstrucción dependiente del objetivo. La
procedencia prospectiva requiere otra flecha,
\[
 \text{especie}\dashrightarrow\text{extensión}
 \longrightarrow\text{ruta}\longrightarrow\text{registro}
 \longrightarrow\text{firma},
\]
y no se deduce de esta expansión. La conclusión fallaría si los \(a_k\) no
tendieran a cero, si la cota del residuo dejara de cumplirse o si la serie de
errores no fuera sumable.
\end{remark}

\section{Bloques globales de la acción de masa}

La acción de masa utilizará la terna
\[
 \Dfour,\qquad s_{120},\qquad \zCorr=135\Dfour^2.
\]
Sus valores se calculan una sola vez. Las coordenadas
\(k,\nu_{120},\nu_{270}\) de una firma especifican cuántas veces interviene
cada bloque en el exponente; no cambian su valor.

Conviene distinguir dos objetos que comparten el índice \(270\):
\[
 \boxed{\zCorr=135\Dfour^2
 =1.66922699663643\times10^{-6}}
\]
y
\[
 \boxed{\sSpec=q_-^{270}-q_+^{270}
 =2.58145536075096\times10^{-11}.}
\]
El primero es un invariante cuadrático construido a partir de \(\Dfour\); el
segundo es un momento espectral del semigrupo. La coincidencia del subíndice no
autoriza su identificación.

\section{Combinación de orden superior}

La combinación global
\[
 \delta\Dfour^{\rm tower}
 =\frac{s_{180}}{24}-\frac{s_{240}}{54}
 -24s_{360}+\frac76s_{420}
\]
evalúa a
\[
 \delta\Dfour^{\rm tower}
 =3.632051471908759123889070967\times10^{-9}.
\]
La coordenada de comparación utilizada en la selección de los
coeficientes se define como el dato escalar
\[
 \delta\Dfour
 :=3.632051471908972784661641720\times10^{-9}.
\]
Por consiguiente, el residuo de la combinación respecto de esa coordenada es
\[
 \delta\Dfour^{\rm tower}-\delta\Dfour
 =-2.1366077257075\times10^{-22}.
\]
La cantidad pequeña es, pues, una diferencia entre dos objetos y no el valor
de la combinación. Todos los sumandos pertenecen a la misma jerarquía y sus
coeficientes son globales. Como la coordenada de comparación intervino en la
selección de los coeficientes, esta igualdad es reconstructiva. La
combinación no se añade de nuevo a la fórmula electrónica que ya emplea
\(\Dfour\), pues esa sustitución duplicaría la misma corrección.

\endgroup
\begingroup
\let\chapter\section
\let\section\subsection
\let\subsection\subsubsection
\let\subsubsection\paragraph
\let\clearpage\relax
\section{Operador bicapa y jerarquía infinita}
\label{sec:jerarquia-espectral-infinita-rev7}
\index{torres!jerarquía infinita}\index{semigrupo espectral}

Sean
\[
 x=A_{\rm rad},\qquad y=C^*_{\rm rad},\qquad
 R_{\rm ch}^2=I,\qquad
 P_\pm^{\rm ch}=\frac{I\pm R_{\rm ch}}2.
\]
El lector de las dos hojas se reúne en el operador
\[
 T_{\rm ang}=\exp(-xI-yR_{\rm ch})
 =q_+P_+^{\rm ch}+q_-P_-^{\rm ch},
\]
con
\[
 q_-=e^{-x+y},\qquad q_+=e^{-x-y}.
\]
Sus momentos espectrales par e impar son
\[
 c_n=\operatorname{Tr}(T_{\rm ang}^n)=q_-^n+q_+^n,\qquad
 s_n=-\operatorname{Tr}(R_{\rm ch}T_{\rm ang}^n)=q_-^n-q_+^n.
\]
La involución \(C^*\mapsto-C^*\) conserva \(c_n\) e invierte \(s_n\).
Equivalentemente,
\[
 c_n=2e^{-nx}\cosh(ny),\qquad
 s_n=2e^{-nx}\sinh(ny).
\]
Ésta es la convención fijada en el capítulo de canales. Para declarar sin
homonimia la presentación alternativa de las fuentes de torres, ponemos
\[
 R_{\rm tow}:=-R_{\rm ch},\qquad
 P_\pm^{\rm tow}=P_\mp^{\rm ch},
\]
\[
 T_{\rm tow}:=e^{-xI+yR_{\rm tow}}=T_{\rm ang},\qquad
 \operatorname{Tr}(R_{\rm tow}T_{\rm tow}^n)
 =-\operatorname{Tr}(R_{\rm ch}T_{\rm ang}^n)=s_n.
\]
Las dos escrituras describen el mismo lector escalar y no constituyen una
segunda orientación de las torres.

\begin{theorem}[Adición, norma y recurrencia]
\label{thm:recurrencia-bicapa-rev7}
Para \(m,n\ge0\),
\[
 c_{m+n}=\frac{c_mc_n+s_ms_n}{2},\qquad
 s_{m+n}=\frac{s_mc_n+c_ms_n}{2},
\]
\[
 c_n^2-s_n^2=4e^{-2nx}.
\]
Si \(p=q_-q_+=e^{-2x}\), ambas sucesiones satisfacen
\[
 X_{n+2}=c_1X_{n+1}-pX_n.
\]
Además,
\[
 \partial_Ac_n=-\frac{n\pi}{180}c_n,
 \qquad
 \partial_As_n=-\frac{n\pi}{180}s_n,
\]
\[
 \partial_{C^*}c_n=\frac{n\pi}{180}s_n,
 \qquad
 \partial_{C^*}s_n=\frac{n\pi}{180}c_n.
\]
\end{theorem}

\begin{proof}
Las identidades de adición resultan al expandir
\((q_-^m\pm q_+^m)(q_-^n\pm q_+^n)\). La norma es
\[
 (q_-^n+q_+^n)^2-(q_-^n-q_+^n)^2
 =4(q_-q_+)^n.
\]
La última igualdad es la recurrencia determinada por las raíces
características \(q_-\) y \(q_+\). Las ecuaciones diferenciales se obtienen
al derivar las expresiones hiperbólicas; el factor \(\pi/180\) pertenece a
la carta angular y no a una selección de especie.
\end{proof}

\begin{theorem}[Semigrupo completo de alturas]
\label{thm:semigrupo-alturas-rev7}
Las alturas de la jerarquía transversal forman
\[
 \mathfrak S
 =\langle90,120\rangle
 =\{90a+120b:a,b\in\mathbb Z_{\ge0},\ a+b>0\}.
\]
Equivalentemente,
\[
 \boxed{\mathfrak S=\{90,120\}\cup\{30r:r\ge6\}.}
\]
\end{theorem}

\begin{proof}
Dividiendo por \(30\), se obtiene el semigrupo
\(\langle3,4\rangle\). Sus elementos positivos son
\(\{3,4\}\cup\{r:r\ge6\}\): \(5\) es la única laguna positiva posterior a
los generadores, y todo \(r\ge6\) se obtiene por inducción sumando \(3\) o
\(4\). Multiplicar de nuevo por \(30\) da la fórmula.
\end{proof}

La presentación conserva también la genealogía de una altura:
\[
 \mathfrak S\cong
 \bigl(\mathbb Z_{\ge0}^2\setminus\{(0,0)\}\bigr)
 \big/\bigl((a+4,b)\sim(a,b+3)\bigr),
\]
pues \(4\cdot90=3\cdot120=360\). La altura no determina por sí sola la
composición que la produjo; las genealogías \(4\cdot90\) y \(3\cdot120\)
permanecen distintas en el registro cronológico enriquecido aunque evalúen
el mismo índice.

Para hacer explícita la recurrencia muestreada en el paso común de treinta,
definamos
\[
 a_{30}:=q_-^{30},\qquad b_{30}:=q_+^{30},\qquad
 u_r:=s_{30r}=a_{30}^r-b_{30}^r.
\]
Entonces
\[
 \boxed{u_{r+2}=(a_{30}+b_{30})u_{r+1}-a_{30}b_{30}u_r.}
\]
El término auxiliar \(u_5=s_{150}\) puede intervenir al propagar la
recurrencia, pero \(150\notin\mathfrak S\); una identidad de cálculo no crea
una nueva altura admisible.

En particular,
\[
90,120,180,210,240,270,300,330,360,390,420,450,480,\ldots
\]
son las primeras alturas. Las ocho alturas destacadas en exposiciones
anteriores son primeras lecturas de \(\mathfrak S\), no su definición ni
su límite.

Como guía humana, su genealogía inicial puede leerse sin convertirla en el
dominio completo:
\[
\begin{array}{c@{\quad}c@{\quad}l}
90&3\cdot30&\text{generador mínimo}\\
120&4\cdot30&\text{generador mínimo}\\
180&2\cdot90&\text{primer modo par}\\
210&90+120&\text{modo compuesto}\\
240&2\cdot120&\text{segundo modo par}\\
270&3\cdot90&\text{repetición de la rama 90}\\
360&3\cdot120&\text{repetición de la rama 120}\\
420&2(90+120)&\text{retorno compuesto}
\end{array}
\]
La última columna registra procedencia estructural; no introduce parámetros
ni afirma que esas ocho filas agoten la jerarquía. Los mismos enteros
\(90,120\) son cardinalidades de incidencia en \(\mathcal F_6\) y
\(\mathcal F_8\), pero esa coincidencia no identifica el operador bicapa con
la inclusión \(\iota_{6,8}\).

La interfaz de los dos generadores es sólo el par
\[
 \mathcal K_{\rm ang}^{90,120}(q_+,q_-):=(s_{90},s_{120}).
\]
No debe confundirse con las resolventes de dos subcolas ni con el funcional
posterior \(K_{6\to8}=12\Delta S_{90}-\Delta S_{120}\), cuya definición y
prueba residen en el \cref{mdg:chap:barbero}. Allí las subcolas viven en la
completación analítica de la jerarquía; no son nuevos generadores ni
coeficientes de especie.

\begin{definition}[Módulo convergente de refinamientos espectrales]
\label{def:modulo-refinamientos-torre-rev7}
El dominio integral de la elevación global es
\[
 \mathcal E_{\mathfrak S}
 :=\left\{\eta\in\mathbb Z^{\mathfrak S}:
 \sum_{h\in\mathfrak S}|\eta_hs_h|<\infty\right\},
 \qquad s_h=q_-^h-q_+^h.
\]
Cada \(\eta\in\mathcal E_{\mathfrak S}\) determina el carácter espectral
\[
 \mathcal R_\eta
 :=\exp\!\left(\sum_{h\in\mathfrak S}\eta_hs_h\right).
\]
\end{definition}

La ley de masas no reside en esta definición: su firma central
\((n_A,s_C,k_{\Delta_4},b_{120},b_\zeta)\in\mathbb Z^5\) se construye antes
desde el registro. La elevación por \(\eta\) actúa después y conserva la
procedencia del término básico \(b_{120}s_{120}\). En particular,
\(N_{120}=b_{120}+\eta_{120}\) es el coeficiente total evaluado, no una
nueva coordenada de firma ni una descomposición inversa única.

La completitud del semigrupo no selecciona por sí sola una especie. Cuando
\(\eta\) se obtiene desde un residual metrológico suministrado, se tiene una
reconstrucción calibrada exacta; la fuerza prospectiva requiere que la ruta,
el acarreo y el devanado produzcan \(\eta\) antes de abrir ese residual. La
recurrencia anterior genera todos los momentos una vez fijado el
refinamiento, sin introducir una ley distinta en cada altura.

\endgroup
\begingroup
\let\chapter\section
\let\section\subsection
\let\subsection\subsubsection
\let\subsubsection\paragraph
\let\clearpage\relax
\section{Par de lectores genealógicos sobre el atlas completo}
\label{sec:operadores-masa-two-way-rev11}
\index{masa!operadores bidireccionales}
\index{atlas!lectores cronológico y dodecafásico}

La ruta enriquecida posee dos representaciones internas complementarias,
\[
 \Gamma\longmapsto
 \bigl(\gamma_{108}(\Gamma),\tau_{12}(\Gamma)\bigr)
 \in\mathbb F_3^{108}\times\{-1,0,+1\}^{12}.
\]
La primera conserva la historia trítica completa; la segunda conserva el
registro dodecafásico firmado. Cada proyección induce un lector natural y ambos
se evalúan antes de abrir masas, etiquetas PDG o residuales.

\HMTClaim{MD-R-MASS-TWOWAY-READER-PAIR}
\subsection{Par de lectores genealógicos \(\mathbf K(\Gamma)=(K_D(\Gamma),K_\Omega(\Gamma))\): dominio, codominio y evaluaciones escalares}

Con la notación de las
\cref{eq:KD-rev11,eq:kappaD-rev11,eq:kappa-omega-rev11}, definimos
\begin{equation}
 \mathbf K(\Gamma)
 :=\bigl(K_D(\Gamma),K_\Omega(\Gamma)\bigr)
 \in\left(\mathbb Z^2\times\tfrac12\mathbb Z\right)
 \times\mathbb Z^3,
 \label{eq:lector-conjunto-masas-rev11}
\end{equation}
donde
\[
 K_D=\left(D_{27}+D_{36},D_1,\frac{D_4-D_3}{2}\right),
 \qquad
 K_\Omega=(\Omega_{90\mid120},54,P_6).
\]
Las evaluaciones escalares son
\begin{align}
 \kappa_D(\Gamma)
 &=D_{27}+D_{36}-\alpha_{\rm HMT}D_1
 +\frac{\Delta_4}{2}(D_4-D_3),
 \label{eq:kappaD-general-rev11}\\
\kappa_\Omega(\Gamma)
&=\Omega_{90\mid120}(\tau_\Gamma)
-54\alpha_{\rm HMT}+P_6(\tau_\Gamma)\Delta_4.
\label{eq:kappaOmega-general-rev11}
\end{align}
La segunda coordenada exige una distinción causal. En APP, el salto
espectral primordial es \(9-5=4\); el coeficiente local \(2\), la compresión
\(6=51-45\) y el despliegue bidimensional \(54=9\cdot6=459-405\) son
invariantes posteriores y tipados, conforme al
\cref{thm:app-jerarquia-cuatro-dos-seis-cincuentaycuatro}. En el lector
cronológico de desplazamientos, en cambio, la coordenada electrónica es
\(B_D=D_1(\Gamma_e)=54\): cuenta fronteras inmediatas y coincide con la
periodicidad cinemática de orden \(54\) del endomorfismo CT108. El
\(54\) fijo de
\(K_\Omega\) es la coordenada común transferida al hacer compatibles ambos
lectores en la reducción electrónica. Su igualdad numérica con el defecto
global de APP constituye un control cruzado posterior; no convierte al
número \(54\) en la diferencia primordial entre suma y producto.

La igualdad \(4\cdot90=3\cdot120=360\) explica la genealogía común de
los dos lectores. \(K_D\) compara desplazamientos de la palabra trítica
\(\gamma_{108}\);
\(K_\Omega\) compara energías de ventanas en la palabra de doce fases. La
coincidencia de origen no identifica sus cocientes de información.
En el dominio certificado \(\Gamma_{324}^{\rm or}\), la periodicidad de
media vuelta hace pares todos los \(D_k\), de modo que
\((D_4-D_3)/2\in\mathbb Z\) y \(K_D\in\mathbb Z^3\). Fuera de ese dominio,
el tercer componente conserva el tipo \(\tfrac12\mathbb Z\).

\subsection{Naturalidad de los operadores bidireccionales y preservación de tipos}

El lector \(D_k\) es invariante bajo traslación temporal, reversión y toda
permutación biyectiva del alfabeto trítico. En las rutas declaradas, la
inversión simultánea de las dos orientaciones actúa por una traslación de
27 pasos discretos y conserva \(K_D\). El lector \(K_\Omega\) es invariante bajo
rotación y reversión dodecafásicas y bajo
\(C_M(\tau)=-\operatorname{rev}(\tau)\). Su forma cuadrática
\[
 \Omega(\tau)=\langle\tau,
 (4W_3^*W_3-3W_4^*W_4)\tau\rangle
\]
posee polarización bilineal entera y satisface la identidad de 2-cociclo.

Los contadores globales de giro, cambio de hoja, repetición de fase, cambio de
prefijo y retorno son constantes sobre las 324 rutas. Por consiguiente, todo
funcional que factorice exclusivamente por esos seis contadores es constante
y cancela en razones. Los lectores de
\eqref{eq:kappaD-general-rev11} y
\eqref{eq:kappaOmega-general-rev11} actúan sobre la historia local y superan
ese control de no factorización.

\begin{ablation}[Necesidad de las dos hojas en la ruta electrónica]
\label{abl:hojas-lector-two-way-rev11}
El lector cronológico produce
\[
 (A_D,B_D,C_D)=
 \begin{cases}
 (80,54,6),&\substack{\text{palabra cronológica conjunta de ambas hojas;}\\
 B_D=D_1(\Gamma_e)=54},\\
 (80,84,-8),&\text{proyección aditiva},\\
 (80,60,13),&\text{proyección multiplicativa}.
 \end{cases}
\]
La terna electrónica completa requiere la historia conjunta de las dos hojas
aritméticas.
\end{ablation}

\subsection{Censos exactos sobre las 324 rutas}

El lector cronológico produce catorce perfiles. La tabla siguiente registra su
censo completo; \(A_D=D_{27}+D_{36}\), \(B_D=D_1\) y
\(C_D=(D_4-D_3)/2\).

\begin{longtable}{@{}rrr r@{}}
\caption{Censo exacto de los catorce perfiles del lector cronológico sobre las
324 rutas orientadas.}
\label{tab:censo-perfiles-kd-rev11}\\
\toprule
\(A_D\)&\(B_D\)&\(C_D\)&multiplicidad\\
\midrule
\endfirsthead
\toprule
\(A_D\)&\(B_D\)&\(C_D\)&multiplicidad\\
\midrule
\endhead
0&24&0&18\\
40&24&2&36\\
40&54&6&18\\
40&78&27&36\\
40&84&30&36\\
60&78&25&18\\
60&84&28&18\\
80&54&6&18\\
80&54&7&18\\
80&66&2&18\\
80&66&3&36\\
80&66&4&18\\
100&66&0&18\\
100&66&2&18\\
\midrule
\multicolumn{3}{r}{total}&324\\
\bottomrule
\end{longtable}

El lector dodecafásico produce nueve perfiles:

\begin{table}[H]
\centering
\caption{Censo exacto de los nueve perfiles del lector dodecafásico sobre
las 324 rutas orientadas.}
\label{tab:censo-perfiles-komega-rev11}
\begin{tabular}{@{}rrr r@{}}
\toprule
\(\Omega\)&\(B_\Omega\)&\(P_6\)&multiplicidad\\
\midrule
80&54&6&198\\
56&54&5&68\\
48&54&5&34\\
36&54&4&8\\
20&54&4&4\\
4&54&2&4\\
40&54&4&4\\
24&54&4&2\\
24&54&2&2\\
\midrule
\multicolumn{3}{r}{total}&324\\
\bottomrule
\end{tabular}
\end{table}

Los lectores coinciden como triples en (18) rutas, siempre con el valor
\((80,54,6)\), y el par \(\mathbf K\) adopta (40) valores. En la sección
canónica \(B1\) de 81 posiciones, \(K_D\) es constante en 41 de las 56
familias y refina el cociente en 77 clases; \(K_\Omega\) es constante en 50
y produce 62 clases; el par produce 78. En el levantamiento de 324
orientaciones, \(K_D\) produce 146 clases y es no constante en todas las
familias; \(K_\Omega\) produce 100 y desciende en 17 familias; el par produce
151 y es no constante en todas. Estos censos prueban que las dos
proyecciones retienen memorias distintas y fijan el dominio de cada
afirmación de descenso.

\subsection{Trazas vectoriales y espectros conjuntos de las \(13\) multisecciones}

En la sección canónica \(B1\), denotamos por
\(\ell_j,\nu_j,d_j,u_j\) las multisecciones de generación \(j\) de los
sectores leptónico cargado, neutrino, quark de tipo abajo y quark de tipo
arriba, respectivamente. Para cada lector
\(r\in\{D,\Omega\}\) y cada componente \(a\in\{1,2,3\}\), definimos el
operador diagonal
\[
 K_{M,a}^r e_\Gamma=K_{r,a}(\Gamma)e_\Gamma,
 \qquad
 \mathbf K_M^r=(K_{M,1}^r,K_{M,2}^r,K_{M,3}^r).
\]
Los tres operadores de \(\mathbf K_M^r\) son autoadjuntos y conmutan. Por
tanto, su traza normalizada se entiende componente a componente,
\[
 \overline{\mathbf K}_M^r
 =\frac{1}{\dim M}
 \bigl(\operatorname{Tr}K_{M,1}^r,
       \operatorname{Tr}K_{M,2}^r,
       \operatorname{Tr}K_{M,3}^r\bigr),
\]
y su espectro es el espectro conjunto, es decir, la distribución con
multiplicidad de los triples simultáneos. Las trece trazas vectoriales son:

\begin{longtable}{@{}c r c c@{}}
\caption{Trazas vectoriales normalizadas de los lectores \(K_D\) y
\(K_\Omega\) en las trece multisecciones canónicas.}
\label{tab:trazas-multisecciones-two-way-rev11}\\
\toprule
\(M\)&\(\dim M\)&\(\overline{\mathbf K}_M^D\)&\(\overline{\mathbf K}_M^\Omega\)\\
\midrule
\endfirsthead
\toprule
\(M\)&\(\dim M\)&\(\overline{\mathbf K}_M^D\)&\(\overline{\mathbf K}_M^\Omega\)\\
\midrule
\endhead
\(\ell_0\)&1&\((80,54,6)\)&\((80,54,6)\)\\
\(\ell_2\)&8&\((65,249/4,17/2)\)&\((80,54,6)\)\\
\(\ell_4\)&16&\((225/4,489/8,21/2)\)&\((141/2,54,11/2)\)\\
\(\nu_1\)&2&\((40,51,29/2)\)&\((60,54,5)\)\\
\(\nu_2\)&4&\((45,63,29/2)\)&\((61,54,5)\)\\
\(\nu_3\)&6&\((170/3,55,34/3)\)&\((206/3,54,11/2)\)\\
\(\nu_4\)&8&\((105/2,237/4,95/8)\)&\((143/2,54,45/8)\)\\
\(d_1\)&2&\((40,51,29/2)\)&\((60,54,5)\)\\
\(d_2\)&4&\((45,63,57/4)\)&\((61,54,5)\)\\
\(d_3\)&6&\((170/3,55,23/2)\)&\((218/3,54,17/3)\)\\
\(d_4\)&8&\((105/2,237/4,49/4)\)&\((149/2,54,23/4)\)\\
\(u_1\)&4&\((65,66,9)\)&\((80,54,6)\)\\
\(u_3\)&12&\((205/3,119/2,113/12)\)&\((74,54,23/4)\)\\
\bottomrule
\end{longtable}

La notación \((a,b,c)^{\times n}\) indica multiplicidad \(n\). El espectro
conjunto completo de \(\mathbf K_D\) es:

\begingroup\footnotesize
\begin{longtable}{@{}c p{.82\textwidth}@{}}
\caption{Espectros conjuntos completos del lector \(K_D\) en las trece
multisecciones canónicas.}
\\
\toprule
\(M\)&\(\operatorname{Spec}_{\rm conj}\mathbf K_M^D\)\\
\midrule
\endhead
\(\ell_0\)&\((80,54,6)^{\times1}\)\\
\(\ell_2\)&\((0,24,0)^{\times1};(40,78,27)^{\times2};(80,54,7)^{\times1};(80,66,2)^{\times1};(80,66,3)^{\times1};(100,66,0)^{\times1};(100,66,2)^{\times1}\)\\
\(\ell_4\)&\((0,24,0)^{\times1}\); \((40,24,2)^{\times2}\); \((40,54,6)^{\times2}\); \((40,84,30)^{\times2}\); \((60,78,25)^{\times3}\); \((80,66,2)^{\times2}\); \((80,66,3)^{\times3}\); \((80,66,4)^{\times1}\)\\
\(\nu_1\)&\((40,24,2)^{\times1};(40,78,27)^{\times1}\)\\
\(\nu_2\)&\((0,24,0)^{\times1};(40,84,30)^{\times1};(60,78,25)^{\times1};(80,66,3)^{\times1}\)\\
\(\nu_3\)&\((40,24,2)^{\times2};(40,78,27)^{\times1};(60,84,28)^{\times1};(80,54,6)^{\times1};(80,66,3)^{\times1}\)\\
\(\nu_4\)&\((0,24,0)^{\times1};(40,24,2)^{\times1};(40,78,27)^{\times2};(40,84,30)^{\times1};(80,54,7)^{\times1};(80,66,2)^{\times1};(100,66,0)^{\times1}\)\\
\(d_1\)&\((40,24,2)^{\times1};(40,78,27)^{\times1}\)\\
\(d_2\)&\((0,24,0)^{\times1};(40,84,30)^{\times1};(60,78,25)^{\times1};(80,66,2)^{\times1}\)\\
\(d_3\)&\((40,24,2)^{\times2};(40,78,27)^{\times1};(60,84,28)^{\times1};(80,54,6)^{\times1};(80,66,4)^{\times1}\)\\
\(d_4\)&\((0,24,0)^{\times1};(40,24,2)^{\times1};(40,78,27)^{\times2};(40,84,30)^{\times1};(80,54,7)^{\times1};(80,66,3)^{\times1};(100,66,2)^{\times1}\)\\
\(u_1\)&\((40,54,6)^{\times1};(60,78,25)^{\times1};(80,66,2)^{\times1};(80,66,3)^{\times1}\)\\
\(u_3\)&\((40,24,2)^{\times2}\); \((40,78,27)^{\times2}\); \((60,84,28)^{\times1}\); \((80,54,6)^{\times3}\); \((80,66,3)^{\times1}\); \((80,66,4)^{\times1}\); \((100,66,0)^{\times1}\); \((100,66,2)^{\times1}\)\\
\bottomrule
\end{longtable}
\label{tab:espectros-kd-multisecciones-rev11}
\endgroup

El espectro conjunto completo de \(\mathbf K_\Omega\) sobre las mismas fibras es:

\begingroup\footnotesize
\begin{longtable}{@{}c p{.82\textwidth}@{}}
\caption{Espectros conjuntos completos del lector \(K_\Omega\) en las trece
multisecciones canónicas.}
\label{tab:espectros-komega-multisecciones-rev11}\\
\toprule
\(M\)&\(\operatorname{Spec}_{\rm conj}\mathbf K_M^\Omega\)\\
\midrule
\endfirsthead
\toprule
\(M\)&\(\operatorname{Spec}_{\rm conj}\mathbf K_M^\Omega\)\\
\midrule
\endhead
\(\ell_0\)&\((80,54,6)^{\times1}\)\\
\(\ell_2\)&\((80,54,6)^{\times8}\)\\
\(\ell_4\)&\((24,54,2)^{\times1};(56,54,5)^{\times4};(80,54,6)^{\times11}\)\\
\(\nu_1\)&\((40,54,4)^{\times1};(80,54,6)^{\times1}\)\\
\(\nu_2\)&\((4,54,2)^{\times1};(80,54,6)^{\times3}\)\\
\(\nu_3\)&\((36,54,4)^{\times1};(56,54,5)^{\times1};(80,54,6)^{\times4}\)\\
\(\nu_4\)&\((36,54,4)^{\times1};(56,54,5)^{\times1};(80,54,6)^{\times6}\)\\
\(d_1\)&\((40,54,4)^{\times1};(80,54,6)^{\times1}\)\\
\(d_2\)&\((4,54,2)^{\times1};(80,54,6)^{\times3}\)\\
\(d_3\)&\((36,54,4)^{\times1};(80,54,6)^{\times5}\)\\
\(d_4\)&\((36,54,4)^{\times1};(80,54,6)^{\times7}\)\\
\(u_1\)&\((80,54,6)^{\times4}\)\\
\(u_3\)&\((56,54,5)^{\times3};(80,54,6)^{\times9}\)\\
\bottomrule
\end{longtable}
\endgroup

Las trazas pueden coincidir entre multisecciones de tipo diferente; el tipo,
la multiplicidad y el espectro completo permanecen como datos del operador.
El certificado reproduce estas tablas fila por fila.

\HMTClaim{MD-R-MASS-FIBRE-OPERATORS}
\subsection{Operadores de fibra y naturalidad del cambio de base}

Distinguimos la fibra canónica \(F_M^{B1}\), formada por las posiciones de
la multisección \(M\) con su orientación publicada, y la fibra orientada
\(F_M^{\rm or}\), formada por sus cuatro levantamientos por posición. Para
\(s\in\{B1,{\rm or}\}\), sea
\(\mathcal H_{M,s}=\bigoplus_{\Gamma\in F_M^s}\mathbb C e_\Gamma\). Definimos
\begin{equation}
 B_{M,s}^r e_\Gamma=\kappa_r(\Gamma)e_\Gamma,
 \qquad
 R_{\rm act}^{B_{M,s}^r}
 =\exp\bigl[(\log R_{\rm act})B_{M,s}^r\bigr],
 \qquad r\in\{D,\Omega\}.
 \label{eq:operadores-fibra-two-way-rev11}
\end{equation}
Cada \(B_{M,s}^r\) es autoadjunto y diagonal en la base de rutas, por lo que
\(R_{\rm act}^{B_{M,s}^r}\) es positivo. Para una masa basal positiva
\(\widehat M_{M,s}^{(0)}\), la forma operatoria condicionada a cada lector es
\begin{equation}
 \widehat M_{M,s}^{\beta,r}
 =R_{\rm act}^{B_{M,s}^r/2}\,
  \widehat M_{M,s}^{(0)}\,
  R_{\rm act}^{B_{M,s}^r/2}.
 \label{eq:ley-operatoria-fibra-rev11}
\end{equation}
Esta expresión es positiva sin imponer conmutatividad. Cuando
\([\widehat M_{M,s}^{(0)},B_{M,s}^r]=0\), reduce a
\(\widehat M_{M,s}^{(0)}R_{\rm act}^{B_{M,s}^r}\).
Un cambio de base unitario \(U\) transforma
\(B_{M,s}^r\mapsto UB_{M,s}^rU^*\) y conserva espectro, traza, determinante y
polinomio característico. La fibra electrónica canónica \(F_e^{B1}\) tiene
rango uno y reduce al escalar del
\cref{thm:reduccion-electron-two-way-rev11}.
Su levantamiento orientado tiene rango cuatro: las rutas \((++)\) y \((--)\)
comparten el perfil electrónico, mientras las orientaciones mixtas conservan
perfiles distintos.

El determinante normalizado
\[
 \operatorname{ndet}(R_{\rm act}^{B_{M,s}^r})
 =R_{\rm act}^{\operatorname{Tr}(B_{M,s}^r)/\dim F_M^s}
\]
es natural bajo permutación de la fibra. Aditividad logarítmica lo distingue
como media geométrica, pero no se usa como sustituto de la realización física.
La salida exacta de refinamiento sigue siendo el par
\((B_{M,s}^D,B_{M,s}^\Omega)\); la sección escalar de una especie se obtiene
por el morfismo de sabor--generación--ruta ya construido en
\cref{eq:morfismo-sabor-ruta-cerrado-rev2,eq:realizacion-fisica-sabor-ruta-cerrada-rev2}:
selecciona la ruta persistente por datos no metrológicos y aplica después el
carácter de su firma. Una funcional escalar arbitraria
\(S(B_{M,s}^D,B_{M,s}^\Omega)\) no define esa selección y permanece sólo como
control de que no se reintroduce la masa objetivo.

\subsection{Escala absoluta, corrección relativa y sensibilidades}

La genealogía de la recta de acción se escribe
\[
 \Sigma_H=\varphi\alpha_{\rm HMT}^{16},
 \qquad \operatorname{ord}_{10}(\Sigma_H)=-34,
\]
\[
 \delta_{2w}=H_5-\eta_{\rm ret}
 =\frac{90}{\pi}\mathscr C_\pi(\alpha_{\rm HMT}),
 \qquad R_{\rm act}=\frac{H_5}{\eta_{\rm ret}},
\]
\[
 \hbar_{\rm pre}=H_5\,10^{-34}\mathcal U_S,
 \qquad
 \hbar_{\rm ret}=\eta_{\rm ret}\,10^{-34}\mathcal U_S.
\]
El factor \(10^{-34}\mathcal U_S\) fija la sección absoluta y cancela en
razones. Para el lector dodecafásico,
\begin{equation}
 \log\!\left[
 \frac{(m_X^{\beta,\Omega}/m_e^{\beta})}
 {(m_X^{(0)}/m_e^{(0)})}\right]
 =\bigl[(\Omega_X-80)+(P_{6,X}-6)\Delta_4\bigr]
 \log R_{\rm act}.
 \label{eq:razon-masa-omega-rev11}
\end{equation}
El término común \(-54\alpha_{\rm HMT}\) cancela también. Para el lector
cronológico,
\begin{equation}
 \log\!\left[
 \frac{(m_X^{\beta,D}/m_e^{\beta})}
 {(m_X^{(0)}/m_e^{(0)})}\right]
 =\bigl[\kappa_D(\Gamma_X)-\kappa_e\bigr]\log R_{\rm act}.
 \label{eq:razon-masa-D-rev11}
\end{equation}
Con \(\log R_{\rm act}=6.932817628081925\ldots\times10^{-10}\),
\[
 \frac{\partial\log m}{\partial\Omega}=6.9328\times10^{-10},
 \quad
 \frac{\partial\log m}{\partial P_6}
 =\Delta_4\log R_{\rm act}=7.7090\times10^{-14},
\]
\[
 \frac{\partial\log m}{\partial\alpha}
 =-54\log R_{\rm act}=-3.74372\times10^{-8},
 \quad
 \left.\frac{\partial\log m}{\partial\Delta_4}\right|_e
 =6\log R_{\rm act}=4.15969\times10^{-9}.
\]
Estas derivadas son sensibilidades del modelo y permanecen separadas de la
incertidumbre experimental de una masa de contraste.

\subsection{Emparejamiento físico y sector nulo}

El emparejamiento tipado sigue la cadena
\[
 \begin{aligned}
 &\text{registro físico tipado}
 \longrightarrow\text{multisección o representación}
 \longrightarrow\text{fibra de rutas}\longrightarrow\text{registro}\\
 &\longrightarrow\text{firma y carácter}
 \longrightarrow(B_{M,s}^D,B_{M,s}^\Omega)
 \longrightarrow\text{carta dimensional}
 \longrightarrow\text{contraste externo}.
 \end{aligned}
\]
El inventario físico tipado distingue registros por objeto, representación y codominio; las 324 rutas son
el motor interno; las 471 filas del catálogo PDG son propiedades externas.
Estos cardinales pertenecen a dominios diferentes. La flecha se construye
desde descriptores físicos no metrológicos, la multisección y la ruta; una
búsqueda por proximidad a una masa pertenece al régimen de calibración.

El sector nulo se bifurca antes del carácter positivo. Para fotón y gluones,
\[
 \kappa_{\rm null}=\mathrm{N/A},
 \qquad m_{\rm null}=0.
\]
La expresión \(R_{\rm act}^{0}=1\) es la unidad del carácter positivo. Los
compuestos concatenan
registros antes de evaluar la firma; los sectores topológicos conservan fusión y
trenza; una sección virtual se clasifica por su horizonte finito de
prolongación.

\subsection{Independencia del constructor de rutas respecto de las masas y clausura del operador}

La verificación estructural utiliza exclusivamente el censo de rutas, el
registro \(81\times108\) y el transductor de las 81 posiciones. Comprueba su
acoplamiento fila a fila, la cobertura
\(81/56/13/324\), las naturalidades, el triple electrónico \((80,54,6)\), los
censos (14) y (9), la coincidencia (18/324), los
(40) pares, descensos diferenciados en \(B1\) y en las 324 orientaciones,
periodicidad 54, inversión por desplazamiento 27 y ablaciones de hoja. El
sector nulo pertenece al dominio exterior al carácter positivo. El cálculo se
realiza antes de introducir masas, CODATA, PDG, residuos metrológicos o
etiquetas nominales.

Este certificado no evalúa \(\alpha_{\rm HMT}\), \(\Delta_4\),
\(R_{\rm act}\), \(\kappa_e\) ni la carta MeV. Esas evaluaciones pertenecen a
la cadena aritmética de la recta de acción y a la realización energética del
\cref{thm:reduccion-electron-two-way-rev11}; esta cadena aritmética permanece
separada de la comprobación combinatoria de los lectores.

La comparación metrológica aplica después los operadores al inventario tipado,
declara unidad, esquema, escala, fecha e incertidumbre, y ejecuta permutaciones
y validación por retención de multisecciones completas. La salida física queda
gobernada por el dominio \(81/56/13/324\), por los dos lectores de ruta y por
los operadores de fibra.

\begin{statusbox}[title={Dominio exacto y realización física cerrada}]
El atlas, la cadena \(C_{81}\to\mathcal L_{108}\) del registro cronológico
enriquecido de \(108\) iteraciones,
la recta de acción,
los lectores \(K_D,K_\Omega\), su reducción electrónica y los operadores de
fibra son resultados internos exactos. El par
\(\mathbf K=(K_D,K_\Omega)\) es la salida completa sobre cada ruta y
determina los dos operadores diagonales de toda fibra. En la fibra
electrónica ambos lectores coinciden. En las fibras no centrales, el par y su
espectro preservan el refinamiento completo, mientras
\(\mathcal S_{\rm phys}\) liga rama, generación, firma fina y, cuando
corresponde, órbita de color con la ruta persistente. El carácter de esa ruta
publica la coordenada escalar antes del contraste metrológico. Por ello no hay
una disyunción entre operador cerrado y masa pendiente: son dos lecturas
tipadas de la misma genealogía. Ambos lectores son co-rectores, y ninguna
proximidad a una masa externa interviene en la selección.
\end{statusbox}
\HMTClaim{MD-R-MASS-PHYSICAL-SELECTION-CLOSED}

\endgroup
\begingroup
\let\chapter\section
\let\section\subsection
\let\subsection\subsubsection
\let\subsubsection\paragraph
\let\clearpage\relax
% Vista editorial exacta de corpus/source/masas/04_persistencia.tex, líneas 365--590. Las líneas 1--364 ya residen exactamente en 73_rev2_electron_lectores_torres.tex.
\subsection{Persistencia del estado enriquecido bajo la extensión superviviente y conservación de la memoria de ruta}
Sea \(x\) un estado APP, \(\tau\) su palabra trítica y \(B_K\) el bloque de memoria al nivel \(K\). El estado relevante para masa tiene la forma

\[
X_K=(x_K,\tau_K,\varepsilon_K,h_K,v_K,c_K,B_K,g_K),
\]
donde \(\varepsilon_K\) es la orientación de hoja, \((h_K,v_K)\) la orientación celular, \(c_K\) el acarreo y \(g_K\) la fase nonádica. Una extensión es superviviente cuando forma una sección compatible de los cilindros sucesivos. Si \(C_{K,j}\), \(0\le j<729\), son los hijos del cilindro en el nivel siguiente, la regla de poda satisface

\[
\sum_{j=0}^{728}G_9(C_{K,j})=1.
\]
El índice único \(j_K\) produce la transición

\[
P_{K+1}=729P_K+j_K.
\]
Al avanzar nueve niveles se conserva la fase,

\[
g(K+9)=g(K),
\]
pero cambian el bloque, el prefijo y la memoria. Por eso el retorno de fase no es retorno de estado. Una partícula estable se define por la compatibilidad de la sección a través de estos retornos, no por la permanencia inmóvil de una coordenada visible.

La ruta observable \(\gamma_{\rm obs}\) es la sombra cronológica de la extensión. El punto que recorre un diagrama es un cursor de lectura; la partícula es la clase completa de persistencia. Esta distinción permite interpretar correctamente el movimiento cuántico: no se exige una bolita siguiendo una trayectoria clásica, sino una historia coherente que admite varias proyecciones compatibles.

\Needspace{7\baselineskip}
\subsection{Registro dodecafásico y firma}
La rueda temporal contiene doce acontecimientos orientados. Su descomposición causal es

\[
12=1+5+6:
\]
un acontecimiento fija el ancla, cinco determinan la incidencia visible y seis forman los pares opuestos que sobreviven al cociente orientado. Si \(s_C^{\rm tot}\) es la suma contraangular entera sobre los doce acontecimientos, el registro conserva

\[
s_C:=s_C^{\rm tot}\in\mathbb Z,
\qquad
s_C^{(12)}:=\frac{s_C^{\rm tot}}6.
\]
El primer objeto es la coordenada entera de la firma; el segundo es su lectura angular dodecafásica. El divisor no es una normalización introducida desde una unidad externa. Es el índice que aparece al reconstruir los seis pares opuestos. La matriz de incidencia dodecafásica posee forma normal de Smith con factores \((1,12)\); el carácter queda bien definido sobre el cociente porque el residuo de orientación se mide en esa estructura.

Ésta es la convención rectora de toda la obra. Las expresiones históricas conservadas que denominaron \(s_C\) a la coordenada ya dividida se leen como \(s_C^{(12)}\); sus cifras y testimonios permanecen intactos, pero ninguna fórmula vigente divide dos veces la misma coordenada.

La pareja de devanados directos es

\[
W_+=6n_A+s_C^{\rm tot},
\qquad
W_-=6n_A-s_C^{\rm tot}.
\]
Como \(s_C=s_C^{\rm tot}\) en la firma publicada, esta fórmula coincide con \(W_\pm=6n_A\pm s_C\). El exponente contraangular usa una sola vez \(s_C^{(12)}=s_C/6\); no vuelve a dividir una coordenada que ya hubiese sido normalizada.

El registro completo publica la firma

\[
\sigma(\Gamma)=
(n_A,s_C,k_{\Delta_4},b_{120},b_\zeta)\in\mathbb Z^5,
\qquad
\zeta_{270}=135\Delta_4^2.
\]
Los cinco componentes tienen tipos distintos. \(n_A\) cuenta el desarrollo angular directo; \(s_C\) registra la suma contraangular entera; \(s_C^{(12)}=s_C/6\) es su lectura angular; \(k_{\Delta_4}\) lee el refinamiento global; \(b_{120}\) marca el canal armónico de altura 120; \(b_\zeta\) pertenece al corrector cuadrático de altura 270. En particular, \(s_{270}\) y \(\zeta_{270}\) no son el mismo objeto: el primero es una traza impar del operador bicapa; el segundo es una coordenada cuadrática del registro.

La palabra electrónica

\[
\tau_{12}=+++---+++---
\]
tiene suma lineal nula, pero su firma par no es nula:

\[
(Q_3,Q_4,Q_6,\rho_0,w)=(-12,-4,12,-6,12).
\]
Tampoco debe identificarse con la firma cruda del ciclo electrónico \((24,0,12,0,-12)\), ni con el origen afín \(v_e=0\). Son cuatro niveles de lectura diferentes: palabra, firma par, firma cruda y origen del carácter relativo.

\Needspace{7\baselineskip}
\subsection{Ángulo, contraángulo y carácter}
Para que la construcción sea autónoma, fijamos aquí las tres coordenadas internas que intervendrán en toda la ley:

\[
\alpha_{\rm HMT}
=\frac{7297352569283800997285105472380663}{10^{36}},
\]
\[
\Delta_4
=\frac{\pi_{\rm HMT}-e_{\rm HMT}\log\pi_{\rm HMT}}{270}
\left(1+\frac{\pi_{\rm HMT}}{729}\right),
\]
y, para \(a>0\) y \(\phi>0\),

\[
H_5(a,\phi)
=\frac12\sqrt{\frac{1000a}{\phi}}-a
+\frac{9a^2}{16}-\frac{5a^3}{9}
+\frac{7a^4}{48}-\frac{a^5}{54}.
\]
En lo sucesivo,

\[
H_5:=H_5(\alpha_{\rm HMT},\varphi_{\rm HMT}).
\]
Estas definiciones no importan una constante metrológica externa: \(\pi_{\rm HMT},e_{\rm HMT},\varphi_{\rm HMT}\) son las evaluaciones coinductivas del lector HMT y \(\alpha_{\rm HMT}\) es la salida racional del cierre dodecafásico. La constante interna \(\alpha_{\rm HMT}\) precede al contraángulo. La rama regional determina primero

\[
A=1000\alpha_{\rm HMT},
\qquad
D_A=\exp\!\left(-\frac{100\pi\alpha_{\rm HMT}}9\right),
\]
\[
\mathcal C_\pi(\alpha_{\rm HMT})
=169\alpha_{\rm HMT}^{6}
+\frac{D_A\alpha_{\rm HMT}^{7}}
       {1-D_A\alpha_{\rm HMT}},
\qquad
y_*=\frac{\pi}{90}(H_5+\alpha_{\rm HMT})
-\mathcal C_\pi(\alpha_{\rm HMT}),
\]
\[
C^*=\frac{180}{\pi}y_*.
\]
Aquí \(A\) y \(C^*\) son coordenadas angulares ya construidas, no magnitudes de acción. Si

\[
\alpha_\angle:=\alpha_{\rm HMT}\,{}^\circ,
\qquad
\eta_{\rm ret}^{(\deg)}:=\frac{C^*}{2}-\alpha_\angle,
\]
entonces

\[
C^*=2\bigl(\eta_{\rm ret}^{(\deg)}+\alpha_\angle\bigr)
\]
es la identidad inversa de retorno, no una segunda regla de selección del contraángulo. Esta escritura evita confundir la mantisa angular \(\eta_{\rm ret}^{(\deg)}\) con la magnitud dimensional \(\hbar_{\rm ret}\). Las cartas en radianes se obtienen mediante la multiplicación por \(\pi/180\). Los canales conjugados son

\[
q_\pm=\exp\left[-\frac\pi{180}(A\pm C^*)\right].
\]
El carácter central de una firma es

\[
\chi_0(\sigma)=
\exp\left(
n_AA_{\rm rad}
+\frac{s_C}{6}C^*_{\rm rad}
+k_{\Delta_4}\Delta_4
+b_\zeta\zeta_{270}
\right).
\]
La coordenada \(b_{120}\) permanece en la firma, pero su contribución se incorpora en la familia de torre. Esta convención impide contar dos veces la misma altura.

El carácter es adimensional. Su función no consiste en convertirse en una constante de acción, sino en actuar por homotecias positivas sobre una recta de masa ya tipada. Para un par de estados \(X,Y\),

\[
\frac{m_Y^{(0)}}{m_X^{(0)}}
=\chi_{\rm full}(\sigma_Y-\sigma_X,\eta_Y-\eta_X).
\]
Esta fórmula separa las razones basales de masa de la elección común de carta absoluta. El cociente beta completo incorpora después \(R_{\rm act}^{\kappa_Y-\kappa_X}\).

\Needspace{7\baselineskip}
\subsection{Semigrupo global de torres de masa y expansión armónica convergente}
Sea \(R^2=I\), con proyectores

\[
P_\pm=\frac{I\pm R}{2},
\qquad
x=\frac{\pi A}{180},
\qquad
y=\frac{\pi C^*}{180},
\]
y

\[
T=e^{-xI-yR}=q_+P_++q_-P_-.
\]
Las trazas par e impar son

\[
c_n=\operatorname{Tr}(T^n)=q_-^n+q_+^n
=2e^{-nx}\cosh(ny),
\]
\[
s_n=-\operatorname{Tr}(RT^n)=q_-^n-q_+^n
=2e^{-nx}\sinh(ny).
\]
Satisfacen

\[
c_n^2-s_n^2=4e^{-2nx},
\]
\[
c_{m+n}=\frac{c_mc_n+s_ms_n}{2},
\qquad
s_{m+n}=\frac{s_mc_n+c_ms_n}{2},
\]
y la recurrencia lineal de orden dos

\[
X_{n+2}=c_1X_{n+1}-e^{-2x}X_n,
\qquad X\in\{c,s\},
\qquad c_0=2,\quad s_0=0.
\]
El conjunto de alturas no se reduce a ocho correcciones visibles. Es el semigrupo

\[
\langle90,120\rangle
=\{90,120\}\cup\{30r:r\ge6\}.
\]
Las alturas \(90,120,180,210,240,270,360,420\) son los primeros testigos genealógicos útiles, no el dominio entero. El carácter refinado es

\[
\log\chi_{\rm full}(\sigma,\eta)
=\log\chi_0(\sigma)
+\sum_{h\in\langle90,120\rangle}N_h(\sigma,\eta)s_h,
\]
con

\[
N_h(\sigma,\eta)=\eta_h+b_{120}\delta_{h,120}.
\]
Así \(s_{120}\) aparece una sola vez, y en altura 270 se conservan separados \(N_{270}s_{270}\) y \(b_\zeta\zeta_{270}\). La convergencia absoluta exige

\[
\sum_h |N_hs_h|<\infty.
\]
La completitud del semigrupo garantiza capacidad representacional de razones positivas; no elige por sí sola los coeficientes físicos de una especie. El regla prospectiva de selección, cuando interviene, debe proceder de ruta, registro genealógico, acarreo, memoria y devanado antes de abrir la masa externa.

\Needspace{7\baselineskip}

\endgroup
\begingroup
\let\chapter\section
\let\section\subsection
\let\subsection\subsubsection
\let\subsubsection\paragraph
\let\clearpage\relax
% Vista editorial exacta de corpus/source/masas/06_operadores.tex, líneas 50--113. Las líneas 1--29 y 30--49 ya residen exactamente en 73_rev2_electron_lectores_torres.tex.
\subsection{Lectores bidireccionales \(K_D\) y \(K_\Omega\): perfiles, coincidencias y operadores por multisección}
La ruta orientada produce dos familias de invariantes. Para los defectos temporales \(D_j\) y la memoria de retorno \(\Omega\), definimos

\[
K_D=\left(D_{27}+D_{36},D_1,\frac{D_4-D_3}{2}\right),
\qquad
K_\Omega=(\Omega,54,P_6).
\]
Sobre las 324 rutas aparecen 14 perfiles distintos de \(K_D\), 9 perfiles de \(K_\Omega\), 40 pares conjuntos y 18 coincidencias. El único perfil común de la ruta electrónica es

\[
K_D(\Gamma_e)=K_\Omega(\Gamma_e)=(80,54,6).
\]
Estos lectores no son etiquetas decorativas. \(K_D\) mide discrepancias de la palabra; \(K_\Omega\) mide el retorno enriquecido. Su igualdad en el centro certifica la clausura bidireccional de rango uno. Fuera del centro producen operadores de fibra cuyos espectros completos deben publicarse aun cuando una realización física posterior seleccione una sección escalar.

\Needspace{7\baselineskip}
\subsection{Refinamiento armónico propio de cada estado}
El refinamiento completo vive en el espacio de coeficientes

\[
\mathcal E_{90,120}
=\left\{\eta=(\eta_h)_{h\in\langle90,120\rangle}:
\sum_h|\eta_hs_h|<\infty\right\}.
\]
La aplicación buscada no recibe una masa. Recibe la ruta enriquecida:

\[
\mathfrak T:\Gamma_{\rm enr}\longrightarrow\mathcal E_{90,120},
\qquad
\Gamma\longmapsto\eta(\Gamma).
\]
Sus coordenadas deben ser funciones explícitas de retornos, acarreo, devanado, memoria y frontera. La condición de naturalidad es

\[
\eta(r_{n,m}\Gamma)=r^{\mathcal E}_{n,m}\eta(\Gamma),
\]
y la convergencia debe acompañarse de una cota de cola calculable. El carácter completo actúa entonces por

\[
\widehat M_F
=\sum_{\gamma\in F}
m_{\rm clk}\,
\chi_{\rm full}(\sigma(\gamma),\eta(\gamma))
|\gamma\rangle\langle\gamma|.
\]
La capacidad de reconstruir una diferencia numérica dada demuestra completitud representativa de la jerarquía; no demuestra por sí sola que \(\mathfrak T\) haya elegido sus coeficientes. La diferencia entre ambas afirmaciones queda preservada en cada fila cuantitativa.

\Needspace{7\baselineskip}
\subsection{Escalarización por acoplamiento}
Sea \(P_{X,a}\) el proyector determinado por el estado \(X\) y el aparato o canal físico \(a\). La distribución de masa es

\[
\mu_{X,a}(B)
=\langle\Psi_X,
P_{X,a}E_{\widehat M_F}(B)P_{X,a}\Psi_X\rangle.
\]
Hay una línea escalar cuando la compresión satisface

\[
P_{X,a}\widehat M_FP_{X,a}=m_XP_{X,a}.
\]
Éste es el significado preciso de «medir es acoplarse». La cifra no estaba oculta en una celda nominal; aparece como autovalor de una compresión física de un operador ya construido. Si la compresión conserva varios autovalores, la predicción correcta es un multiplete o una distribución, no una media elegida para coincidir con una referencia.

\Needspace{7\baselineskip}


\endgroup
\begingroup
\let\chapter\section
\let\section\subsection
\let\subsection\subsubsection
\let\subsubsection\paragraph
\let\clearpage\relax
\subsection{Selección física como relación natural}
Fuera del centro electrónico, una especie no debe elegir por decreto una única celda. Las doce multisecciones no centrales son fibras estables bajo la inversión celular y no admiten, en general, una sección puntual equivariante. La realización correcta empieza por una relación

\[
\mathscr S\subset \mathcal P\times\mathcal R_{324},
\qquad
\mathscr S(X)=\{\gamma:(X,\gamma)\in\mathscr S\},
\]
cuyo valor puede ser un punto, una órbita o una fibra. Para el electrón, \(\mathscr S(e)=\{\gamma_e\}\) es de rango uno porque \((5,5)\) es el único punto fijo. Para color, sabor y generación, \(\mathscr S(X)\) conserva la órbita correspondiente hasta que un acoplamiento físico la comprime.

La relación es admisible si satisface simultáneamente:

\begin{enumerate}
\item depende sólo de representación, orientación, memoria, frontera, carga,
sabor, color y generación construidos antes del contraste;

\item conmuta con conjugación de partícula y antipartícula;
\item es equivariante respecto de la acción cromática y de la inversión celular;
\item es compatible con refinamiento, retorno nonádico y lectura dodecafásica;
\item conserva toda multiplicidad que no haya sido eliminada por un acoplamiento;
\item permanece invariante al permutar o reemplazar los valores externos de masa.
\end{enumerate}
El sexto punto es el control decisivo. Si alterar la columna experimental cambia \(\mathscr S(X)\), la construcción es retrospectiva. Si la relación permanece y cambia sólo el residual de contraste, la dirección causal está preservada.

\Needspace{7\baselineskip}

\endgroup

\section{Morfismo de escalarización y obstrucciones de la realización dimensional}
Fuera de la fibra electrónica de rango uno, el espectro no selecciona por sí
mismo un autovalor físico. Un escalar requiere un proyector, estado o funcional
de acoplamiento declarado antes del contraste externo.
